% !TeX root = ../../Book3.tex
\chapter{全纯函数与幂级数}
\label{b3:ch:01}
\BookChapterNumber{b3:ch:01}
% 本章摘要（不计入编号）
\begin{chapterabstract}
复导数要求差商从所有方向趋近同一个值。这个局部条件最终会迫使函数具有收敛的幂级数展开，带来远强于实可微的刚性。本章从复差商与 Cauchy--Riemann 方程开始，证明幂级数的收敛与微分性质，再借助圆盘内的积分公式建立全纯与解析的等价性。随后构造指数、对数、三角函数和 Möbius 变换，最后由零点分解得到恒等定理、局部逆函数定理与开映射定理。
\end{chapterabstract}

% 本章目标（不计入编号）
\begin{learninggoals}
\begin{itemize}
\item\label{b3:item:01-goal-derivative}能区分复可微、实可微与全纯，正确使用 Cauchy--Riemann 方程，并解释非零复导数的几何意义。
\item\label{b3:item:01-goal-series}能计算收敛半径与闭子圆盘上的截断误差，判断逐项微分的范围，并说明局部积分公式如何推出幂级数展开。
\item\label{b3:item:01-goal-elementary}能在明确的定义域和分支约定下使用初等复函数，核对 Möbius 变换在有限点与无穷远点的行为。
\item\label{b3:item:01-goal-local}能利用零点阶数和连通性判断函数的唯一性，构造局部标准坐标并计算附近原像的数目，区分局部可逆、全局单射与开映射。
\end{itemize}
\end{learninggoals}

实函数可以无限次可微，却不由一点的全部导数确定。复函数会发生什么变化？先比较 \(f(z)=z^2\) 与 \(g(z)=\overline z\)：前者的差商沿所有方向有同一极限，后者沿实方向和虚方向给出相反的值。要求所有方向相容，看似只增加一个微分条件，却会迫使函数具有收敛幂级数，并使它在一小段曲线上的值决定整个连通区域中的值。这条推理需要两座桥：先用积分控制微分所不能直接控制的余项，再用幂级数的首个非零项描述局部行为。实线性映射的分解、几何级数的尾项以及局部开根的计算，将分别把这些步骤具体化。

\ProseParagraphBreak

比较实微分与复微分，有助于理解这种刚性：一般实线性映射可以分别改变两个方向，复线性映射却只能同时旋转和伸缩。局部求逆时，第一卷的实逆函数定理也会在这个二维图景中发挥作用。

\ProseParagraphBreak

\subsecref{b3:subsec:010204}先证明三角形上的 Goursat 定理与圆盘内的 Cauchy 积分公式，再由此推出全纯函数的局部幂级数展开。\chapref{b3:ch:02}将讨论一般曲线与定义域的拓扑条件。熟悉这些结果的读者仍需留意本卷关于对数主值和保向共形的约定。

\ProseParagraphBreak

John B. Conway 的《Functions of One Complex Variable I》\cite[Chapter III, \S1, pp. 30--31]{Conway1978Functions}从绝对收敛和几何级数出发讨论收敛半径，可与\secref{b3:sec:0102}的证明比较。Theodore W. Gamelin 的《Complex Analysis》\cite{Gamelin2001Complex}则依次讲述初等函数、解析函数、复积分和幂级数，与本章先建立幂级数性质、再构造初等函数的安排不同。

% 本章来源实读范围见参考文献原文件/来源记录.md；证明义务以现行正文和工作记录/审核进度.md为准。
% Conway 公开样章实读印刷30--31页；Gamelin仅核验官方书目与目录。
% 跨卷调用稳定标签与当前编号登记于跨卷引用表.md。

% !TeX root = ../../Book3.tex
\section{复可微与全纯}
\label{b3:sec:0101}

% 本节从差商与实全微分独立证明基本法则，不调用跨卷编号。
% 反例与旋转伸缩图为按本节辨析目标自行组织的表达。
% 未使用全纯函数导数连续、局部幂级数、开映射或局部逆函数结论。

把复平面识别为实平面之后，复函数可以看成两个实变量的向量值函数。复导数的存在比两个偏导数的存在要求更强：无论增量沿哪个方向趋于零，差商都必须趋于同一个复数。本节用实微分刻画这一条件，并说明非零复导数为何对应于局部的旋转与伸缩。

\subsection{复导数}
\label{b3:subsec:010101}

以下写 \(z=x+\mathrm{i}y\)，其中 \(\mathrm{i}^2=-1\)，并用 \(\overline z=x-\mathrm{i}y\) 和 \(\lvert z\rvert=(x^2+y^2)^{1/2}\) 分别表示共轭与模。所有邻域和极限均取复平面的欧氏距离。设 \(U\subseteq\mathbb C\) 为开集；这一条件保证每个 \(a\in U\) 附近的任意充分小增量都能参与差商。
\symidx{\(\overline z\)：复数 \(z\) 的共轭}
\symidx{\(\lvert z\rvert\)：复数 \(z\) 的模}

\begin{definition}\label{b3:def:01-complex-derivative}
设 \(f:U\rightarrow\mathbb C\)，且 \(a\in U\)。若有限极限
\[
 f'(a)\coloneqq\lim_{\substack{h\to0\\h\ne0}}
 \frac{f(a+h)-f(a)}{h}
\]
存在，则称 \(f\) 在 \(a\) 处\term[fu-kewei]{复可微}[complex differentiable]，称 \(f'(a)\) 为该点的\term[fu-daoshu]{复导数}[complex derivative]。
\end{definition}
\symidx{\(f'(a)\)：函数在 \(a\) 处的复导数}

极限存在等价于存在复数 \(c\)，使
\begin{equation}\label{b3:eq:01-complex-linear-approximation}
 f(a+h)=f(a)+ch+o(\lvert h\rvert).
\end{equation}
这里 \(o(\lvert h\rvert)\) 表示一个复余项，其模除以 \(\lvert h\rvert\) 趋于零；下文的 \(O(|h|^k)\) 或 \(O(h^k)\) 表示模不超过某个常数倍 \(|h|^k\) 的余项。上述等价性来自 \(\lvert rh^{-1}\rvert=\lvert r\rvert/\lvert h\rvert\)。系数只能是 \(c=f'(a)\)。这一表达立即说明复可微蕴含连续，也说明微分的首项是乘以一个固定复数，而不能是任意实线性变换。

\ProseParagraphBreak

例如 \(f(z)=z^2\) 满足 \(f(a+h)-f(a)=2ah+h^2\)，故处处复可微，且 \(f'(a)=2a\)。共轭函数则完全不同：当 \(h\) 为非零实数时，\(\overline h/h=1\)；当 \(h\) 为非零纯虚数时，\(\overline h/h=-1\)。因此 \(z\mapsto\overline z\) 在任何一点都不复可微，尽管它作为实平面上的映射处处具有连续的各阶偏导数。

\begin{proposition}\label{b3:prop:01-complex-derivative-rules}
设 \(f,g\) 在 \(a\) 复可微，\(\alpha,\beta\in\mathbb C\)。则线性组合与乘积在 \(a\) 复可微，且
\[
 (\alpha f+\beta g)'(a)=\alpha\cdot f'(a)+\beta\cdot g'(a),
 \quad
 (fg)'(a)=f'(a)\cdot g(a)+f(a)\cdot g'(a).
\]
若 \(g(a)\ne0\)，则商在 \(a\) 的某邻域有定义，并且
\[
 \left(\frac fg\right)'(a)
 =\frac{f'(a)\cdot g(a)-f(a)\cdot g'(a)}{g(a)^2}.
\]
\end{proposition}
\begin{proof}
线性组合的差商直接按极限的线性性处理。乘积的增量写成
\[
 f(a+h)\cdot\bigl(g(a+h)-g(a)\bigr)
 +g(a)\cdot\bigl(f(a+h)-f(a)\bigr),
\]
除以 \(h\) 后，利用 \(f\) 在 \(a\) 连续即可取极限。若 \(g(a)\ne0\)，连续性保证充分小时 \(\lvert g(a+h)\rvert\geq\lvert g(a)\rvert/2\)。于是
\[
 \frac{g(a+h)^{-1}-g(a)^{-1}}h
 =-\frac{g(a+h)-g(a)}{h\cdot g(a+h)\cdot g(a)}
 \longrightarrow-\frac{g'(a)}{g(a)^2}.
\]
再应用乘积公式得到商法则。
\end{proof}

复合运算的求导由\term[lianshi-faze]{链式法则}[chain rule]控制。直接把两个差商相乘时，可能遇到内层函数在邻近点取相同值的情形；用线性近似可以同时处理这些点。

\begin{proposition}[链式法则]\label{b3:prop:01-complex-chain-rule}
若 \(f\) 在 \(a\) 复可微，\(F\) 定义在 \(f(a)\) 的开邻域内，并在 \(f(a)\) 复可微，则复合函数在 \(a\) 的某邻域有定义，且
\[
 (F\circ f)'(a)=F'\bigl(f(a)\bigr)\cdot f'(a).
\]
\end{proposition}
\begin{proof}
令 \(b=f(a)\)，并写
\[
 \begin{aligned}
 f(a+h)-b&=f'(a)h+r(h),&r(h)&=o(\lvert h\rvert),\\
 F(b+k)-F(b)&=F'(b)k+s(k),&s(k)&=o(\lvert k\rvert).
 \end{aligned}
\]
把最后一式在 \(k=0\) 时的余项取为零。第一个展开给出 \(\lvert k\rvert=O(\lvert h\rvert)\)，因此代入 \(k=f(a+h)-b\) 后有 \(s(k)=o(\lvert h\rvert)\)，即使某些非零 \(h\) 对应 \(k=0\) 也成立。复合函数的线性首项遂为 \(F'(b)\cdot f'(a)h\)。其邻域定义由 \(f\) 在 \(a\) 的连续性保证。
\end{proof}

\subsection{Cauchy--Riemann 方程}
\label{b3:subsec:010102}

写 \(f=u+\mathrm{i}v\)，把 \(u,v\) 看作两个实变量的实值函数。作为实映射，\(f\) 在 \(a\) 可微是指存在实线性映射 \(A:\mathbb R^2\rightarrow\mathbb R^2\)，使
\[
 f(a+h)-f(a)=Ah+o(\lvert h\rvert).
\]
该映射称为\term[frechet-weifen]{Fréchet 微分}[Fréchet differential]，记作 \(Df(a)\)。若存在，它的矩阵由偏导数确定为
\[
 Df(a)=
 \begin{pmatrix}
 u_x(a)&u_y(a)\\
 v_x(a)&v_y(a)
 \end{pmatrix}.
\]
下文把这一条件简称为“实可微”。复数乘法还要求此矩阵具有特殊形状。
\symidx{\(Df(a)\)：复函数作为实映射的 Fréchet 微分}

\begin{theorem}[Cauchy--Riemann]\label{b3:thm:01-cauchy-riemann-criterion}
函数 \(f=u+\mathrm{i}v\) 在 \(a\) 复可微，当且仅当它在该点实可微，并满足\term[cauchy-riemann-fangcheng]{Cauchy--Riemann 方程}[Cauchy--Riemann equations]
\begin{equation}\label{b3:eq:01-cauchy-riemann}
 u_x(a)=v_y(a),\quad u_y(a)=-v_x(a).
\end{equation}
成立时，
\[
 f'(a)=u_x(a)+\mathrm{i}v_x(a)=v_y(a)-\mathrm{i}u_y(a).
\]
\end{theorem}
\begin{proof}
若 \(f'(a)=\alpha+\mathrm{i}\beta\)，由\eqref{b3:eq:01-complex-linear-approximation}得到实微分
\[
 Df(a)=
 \begin{pmatrix}\alpha&-\beta\\ \beta&\alpha\end{pmatrix}.
\]
比较矩阵元素即得方程及导数公式。反过来，若实微分存在且满足\eqref{b3:eq:01-cauchy-riemann}，令 \(c=u_x(a)+\mathrm{i}v_x(a)\)，则这个实微分恰好是乘法 \(h\mapsto ch\)。实可微的余项因此给出\eqref{b3:eq:01-complex-linear-approximation}，从而复导数存在。
\end{proof}

定理中的实可微不能删去。定义
\[
 f(x+\mathrm{i}y)=
 \begin{cases}
 \dfrac{x^3+\mathrm{i}y^3}{x^2+y^2},&(x,y)\ne(0,0),\\[4pt]
 0,&(x,y)=(0,0).
 \end{cases}
\]
由 \(x^6+y^6\leq(x^2+y^2)^3\) 得 \(\lvert f(z)\rvert\leq\lvert z\rvert\)，所以它在原点连续。沿两坐标轴计算，得到 \(u_x(0)=v_y(0)=1\)、\(u_y(0)=v_x(0)=0\)，满足 Cauchy--Riemann 方程。但当 \(t\ne0\) 时，
\[
 \frac{f(t)}t=1,\quad
 \frac{f(t+\mathrm{i}t)}{t+\mathrm{i}t}=\frac12.
\]
差商没有极限。这个例子同时具有连续性与满足方程的四个偏导数，仍缺少对任意小增量的统一线性近似。

\begin{corollary}\label{b3:cor:01-continuous-partial-criterion}
若 \(u,v\) 的四个一阶偏导数在 \(a\) 的某邻域存在，并在 \(a\) 连续，则 \(f\) 在 \(a\) 复可微当且仅当 Cauchy--Riemann 方程在 \(a\) 成立。
\end{corollary}
\begin{proof}
只须证明实可微性。写 \(a=x+\mathrm{i}y\)。对 \(w=u\) 或 \(w=v\)，把增量从 \((x,y)\) 到 \((x+s,y+t)\) 分成先水平后竖直的两段。一元中值定理给出
\[
 w(x+s,y+t)-w(x,y)
 =s\cdot w_x(\xi,y)+t\cdot w_y(x+s,\eta),
\]
其中 \(\xi\) 介于 \(x\) 与 \(x+s\) 之间，\(\eta\) 介于 \(y\) 与 \(y+t\) 之间；零长度的一段直接记为零。减去 \(s\cdot w_x(a)+t\cdot w_y(a)\)，由两偏导在 \(a\) 连续，余项为 \(o\bigl((s^2+t^2)^{1/2}\bigr)\)。对两个分量合并后应用\cref{b3:thm:01-cauchy-riemann-criterion}。
\end{proof}

这里的连续偏导条件是检验复可微的一条充分途径。直接从差商定义出发时，不需要先假定导数随基点连续。

\subsection{全纯函数}
\label{b3:subsec:010103}

\begin{definition}\label{b3:def:01-holomorphic}
设 \(U\subseteq\mathbb C\) 为开集。若 \(f:U\rightarrow\mathbb C\) 在 \(U\) 的每一点复可微，则称它在 \(U\) 上\term[quan-chun]{全纯}[holomorphic]，称其为\term[quan-chun-hanshu]{全纯函数}[holomorphic function]，全体这样的函数记作 \(\mathcal H(U)\)。若 \(f\) 在某个包含 \(a\) 的开邻域上全纯，则称它在 \(a\) 附近全纯。在整个复平面上全纯的函数称为\term[zheng-hanshu]{整函数}[entire function]。
\end{definition}
\symidx{\(\mathcal H(U)\)：开集 \(U\) 上的全纯函数类}

全纯性对一个开集内的所有点提出要求，强于在一个指定点复可微。例如 \(f(z)=\lvert z\rvert^2\) 在原点的差商为 \(\overline h\)，所以 \(f'(0)=0\)。若 \(a=x+\mathrm{i}y\ne0\)，沿实轴与虚轴增量所得的差商极限分别为 \(2x\) 与 \(-2\mathrm{i}y\)，两者相等只可能发生在 \(a=0\)。故该函数在原点复可微，却不在原点的任何邻域上全纯。

\ProseParagraphBreak

由\cref{b3:prop:01-complex-derivative-rules}，多项式都是整函数，有理函数在分母不为零的开集上全纯；全纯函数的线性组合与乘积仍全纯。\cref{b3:prop:01-complex-chain-rule}还保证定义范围内的复合全纯。商法则只允许在分母不为零处使用。至于某个零点是否能通过重新赋值补上，需要另作极限分析，不能单凭代数表达式下结论。

\ProseParagraphBreak

在本卷中，非空连通开集称为\term[quyu]{区域}[domain]。连通条件尚未出现在全纯函数的定义里，但它让局部结论能够传播。例如，区域上的实值全纯函数只能是常数。

\begin{proposition}\label{b3:prop:01-real-valued-holomorphic}
若 \(U\) 为区域，\(f\in\mathcal H(U)\)，且 \(f(U)\subseteq\mathbb R\)，则 \(f\) 为常数。
\end{proposition}
\begin{proof}
写 \(f=u+\mathrm{i}v\)，则 \(v=0\)。由 Cauchy--Riemann 方程，\(u_x=u_y=0\)。对 \(U\) 中任意足够小的开矩形，沿其水平线段及竖直线段分别使用一元中值定理，得到 \(u\) 在该矩形上恒定。因此 \(u\) 局部恒定。固定 \(a\in U\)，集合 \(E=\{\,z\in U\mid u(z)=u(a)\,\}\) 非空，局部恒定性使 \(E\) 与其在 \(U\) 内的补集都为开集。由 \(U\) 连通，\(E=U\)。
\end{proof}

全纯函数的定义只含一次复导数。\secref{b3:sec:0102}将证明它还能在每一点附近展开成幂级数，因而具有任意阶导数。这个结论是需要证明的正则性，不能在当前的判别中倒过来充当假设。

\subsection{共形性}
\label{b3:subsec:010104}

若实线性映射 \(A:\mathbb C\rightarrow\mathbb C\) 还满足 \(A(\mathrm{i}h)=\mathrm{i}\cdot A(h)\)，则它与复数数乘相容，称为\term[fu-xianxing-yingshe]{复线性映射}[complex-linear map]。令 \(c=A(1)\)，利用实线性可得
\[
 A(x+\mathrm{i}y)=x\cdot A(1)+y\cdot A(\mathrm{i})=(x+\mathrm{i}y)c.
\]
因此复平面上的复线性映射恰好是乘以一个复数。若 \(c\ne0\)，写 \(c=\rho(\cos\theta+\mathrm{i}\sin\theta)\)，其中 \(\rho=\lvert c\rvert>0\)，则
\[
 A=\rho
 \begin{pmatrix}
 \cos\theta&-\sin\theta\\
 \sin\theta&\cos\theta
 \end{pmatrix}.
\]
这表示先旋转角 \(\theta\)，再把长度都乘以 \(\rho\)。\cref{b3:fig:01-complex-similarity}给出一组向量的变化；两个方向的夹角保持不变。

\begin{figure}[htbp]
\centering
\begin{tikzpicture}[>=stealth,scale=0.95]
 \begin{scope}
  \draw[->,AnalysisGray] (-0.45,0)--(2.05,0);
  \draw[->,AnalysisGray] (0,-0.4)--(0,1.85);
  \fill (0,0) circle (1.4pt);
  \node[below left] at (0,0) {\(0\)};
  \draw[->,thick,AnalysisBlue] (0,0)--(1.45,0) node[below] {\(h_1\)};
  \draw[->,thick,AnalysisBlue] (0,0)--(0.725,1.256) node[above right] {\(h_2\)};
  \draw[AnalysisBlue] (0.45,0) arc[start angle=0,end angle=60,radius=0.45];
  \node at (0.71,0.38) {\(\alpha\)};
 \end{scope}
 \draw[->,thick] (2.25,0.8)--(4.10,0.8)
  node[midway,above] {\(h\mapsto ch\)};
 \begin{scope}[xshift=5.2cm]
  \draw[->,AnalysisGray] (-0.55,0)--(2.0,0);
  \draw[->,AnalysisGray] (0,-0.4)--(0,2.1);
  \fill (0,0) circle (1.4pt);
  \node[below left] at (0,0) {\(0\)};
  \draw[->,thick,AnalysisBlue] (0,0)--(1.575,0.909) node[right] {\(ch_1\)};
  \draw[->,thick,AnalysisBlue] (0,0)--(0,1.819) node[above right] {\(ch_2\)};
  \draw[AnalysisBlue] (0.433,0.25) arc[start angle=30,end angle=90,radius=0.5];
  \node at (0.41,0.72) {\(\alpha\)};
  \draw[AnalysisGray] (0.85,0) arc[start angle=0,end angle=30,radius=0.85];
  \node at (1.12,0.25) {\(\theta\)};
 \end{scope}
\end{tikzpicture}
\caption[复线性微分的旋转与伸缩]{复线性微分的旋转与伸缩。图中取 \(\theta=\pi/6\)、\(\rho>1\)，相交方向的有向夹角 \(\alpha\) 不变。}
\label{b3:fig:01-complex-similarity}
\end{figure}

\begin{definition}\label{b3:def:01-conformal}
设 \(f\) 在 \(a\) 实可微。若 \(Df(a)\) 是非零的旋转伸缩，即具有上述 \(\rho>0\) 的形式，则称 \(f\) 在 \(a\) \term[gongxing]{共形}[conformal]。本卷的这一用语约定保持定向；在开集的每一点共形时，称该映射在此开集上共形。
\end{definition}

对实线性可逆映射，这个定义也等价于保持有向夹角。事实上，标准基 \(1,\mathrm{i}\) 的像必须垂直且定向为正，可写成 \(r e,s\mathrm{i}e\)，其中 \(r,s>0\)、\(\lvert e\rvert=1\)。再考察 \(1\) 与 \(1+\mathrm{i}\) 的夹角 \(\pi/4\)，得到 \(s/r=1\)，故两个伸缩因子相同。反向的角度保持直接由旋转伸缩形式得到。

\begin{proposition}\label{b3:prop:01-conformal-criterion}
函数 \(f\) 在 \(a\) 共形，当且仅当它在 \(a\) 复可微且 \(f'(a)\ne0\)。此时
\[
 \det Df(a)=\lvert f'(a)\rvert^2>0.
\]
若两条实参数曲线在 \(a\) 相交，且在交点具有非零切向量，则它们的像在该点仍具有非零切向量，二者的有向夹角保持不变。
\end{proposition}
\begin{proof}
由\cref{b3:thm:01-cauchy-riemann-criterion}，复可微恰好表示实微分是复数乘法；乘数不为零时，它恰好是非零旋转伸缩。矩阵形式给出行列式等式。

设两条曲线满足 \(\gamma_j(0)=a\)、\(\gamma_j'(0)=v_j\ne0\)。把 \(\gamma_j(t)=a+tv_j+o(\lvert t\rvert)\) 代入\eqref{b3:eq:01-complex-linear-approximation}，得到
\[
 (f\circ\gamma_j)'(0)=f'(a)v_j\ne0.
\]
两个像向量相除所得的比值为 \(v_2/v_1\)，与原向量的比值相同，因而有向夹角相同。
\end{proof}

共轭映射虽保持夹角的大小，却把有向夹角反号，不满足这里的保向约定。导数为零时也不能沿用上述结论：对 \(f(z)=z^2\)，从原点出发的射线辐角会变为原来的两倍，而在原点可微的实参数曲线的像切向量都被微分压成零向量。最后，共形条件本身是局部的微分条件，并未断言整个定义域上的单射性；例如 \(z^2\) 在 \(\mathbb C\setminus\{0\}\) 处处共形，却把 \(z\) 与 \(-z\) 映到同一点。单射性以及局部逆函数的存在将在\subsecref{b3:subsec:010403}中分别讨论。

% !TeX root = ../../Book3.tex
\section{幂级数}
\label{b3:sec:0102}

多项式可以直接求导，而无穷级数的求导必须受到收敛条件约束。幂级数的特殊之处在于：只要把讨论限制在收敛圆的内部，级数本身及它的各阶导数就都有统一控制。更深的一步是反过来证明，每个全纯函数局部都来自这样的级数。这个结论不能仅由 Cauchy--Riemann 方程的形式推出。

\subsection{收敛半径}
\label{b3:subsec:010201}

固定 \(a\in\mathbb C\) 及复数列 \((c_n)_{n\geq0}\)。形如
\[
\sum_{n=0}^\infty c_n(z-a)^n
\]
的级数称为以 \(a\) 为中心的\term[miji-shu]{幂级数}[power series]。常数项为 \(c_0\)，在 \(z=a\) 处级数总是收敛。记 \(D(a,r)=\{z\in\mathbb C:|z-a|<r\}\)，并约定 \(D(a,\infty)=\mathbb C\)。
\symidx{\(D(a,r)\)：中心为 \(a\)、半径为 \(r\) 的开圆盘}

\ProseParagraphBreak

先观察收敛向中心传播的原因。若级数在某个 \(z_1\ne a\) 处收敛，置 \(\rho=|z_1-a|\)，则它的通项有界，存在 \(M_\rho<\infty\) 使
\[
 |c_n|\rho^n\le M_\rho\quad(n\ge0).
\]
对任何 \(0<r<\rho\)，在 \(|z-a|\le r\) 上都有
\[
 |c_n(z-a)^n|\le M_\rho(r/\rho)^n.
\]
右端是可求和的几何数列。因此，一个非中心点的收敛就保证更小闭圆盘上的绝对一致收敛；这个论证没有假定原级数在 \(z_1\) 绝对收敛。控制范围天然呈圆形，是因为各项的绝对值只含距离 \(|z-a|\)。

\ProseParagraphBreak

下面的\term[cauchy-hadamard-gongshi]{Cauchy--Hadamard 公式}[Cauchy--Hadamard formula]从系数确定收敛的范围。

\Needspace{6\baselineskip}
\begin{theorem}[Cauchy--Hadamard]\label{b3:thm:01-power-radius}
设 \(a\in\mathbb C\)、\(c_n\in\mathbb C\ (n\ge0)\)，考虑幂级数 \(\sum_{n=0}^{\infty}c_n(z-a)^n\)。令 \(q=\limsup_{n\to\infty}|c_n|^{1/n}\)，并规定 \(1/0=\infty\)、\(1/\infty=0\)。则 \(R=1/q\) 满足：级数在 \(|z-a|<R\) 时绝对收敛，在 \(|z-a|>R\) 时发散。对任意 \(0<r<R\)，级数在闭圆盘 \(\overline{D(a,r)}\) 上绝对一致收敛。
\end{theorem}
\begin{proof}
若 \(0<r<R\)，可选 \(b>q\) 使 \(br<1\)。由上极限的定义，对充分大的 \(n\) 有 \(|c_n|\leq b^n\)，故在 \(|z-a|\leq r\) 上，级数尾部由几何级数 \(\sum_n(br)^n\) 控制。它既证明绝对收敛，也给出与 \(z\) 无关的尾部误差估计。

若 \(s=|z-a|>R\)，则可选 \(b<q\) 使 \(bs>1\)，并有无穷多个 \(n\) 满足 \(|c_n|^{1/n}>b\)。对这些指标，\(|c_n(z-a)^n|>(bs)^n\)，所以通项不趋于零。上述选择在 \(q=0\) 或 \(q=\infty\) 时仍按相应方向适用：前者任意有限闭圆盘上收敛，后者每个非中心点处发散。
\end{proof}

对复系数幂级数 \(\sum_{n\ge0}c_n(z-a)^n\)，数 \(R=(\limsup_{n\to\infty}|c_n|^{1/n})^{-1}\)（倒数按扩展非负实数约定）称为\term[shoulian-banjing]{收敛半径}[radius of convergence]。定理没有判定边界。半径同为 \(1\) 的三个级数
\[
\sum_{n=0}^\infty z^n,\quad
\sum_{n=1}^\infty\frac{z^n}{n},\quad
\sum_{n=1}^\infty\frac{z^n}{n^2}
\]
分别展示不同情形：第一个在单位圆上处处发散；第二个在 \(z=1\) 处发散、在 \(z=-1\) 处由交错级数判据收敛；第三个在整个闭单位圆盘上绝对一致收敛，因为尾部由 \(\sum n^{-2}\) 控制。因此，不能把开圆盘内的结论直接延到边界。

\ProseParagraphBreak

两个极端值也会实际出现。若 \(c_n=n!\)，则对 \(n\ge2\)，最后至少 \(n/2\) 个因子不小于 \(n/2\)，故
\[
 (n!)^{1/n}\ge(n/2)^{1/2}\longrightarrow\infty.
\]
相应级数的半径为零，只有中心点收敛。若 \(c_n=n^{-n}\)（\(n\ge1\)），则 \(|c_n|^{1/n}=1/n\to0\)，半径为无穷大。这些约定并非只为统一记号；没有正收敛半径的形式级数，不能直接当作邻域中的函数使用。

\ProseParagraphBreak

闭子圆盘上的控制还能给出具体误差。若 \(0<r<\rho<R\)，并取 \(M_\rho=\sup_{n\ge0}|c_n|\rho^n<\infty\)，则对每个整数 \(N\ge0\)，有
\begin{equation}\label{b3:eq:01-power-tail}
 \sup_{|z-a|\le r}\left|f(z)-\sum_{n=0}^{N}c_n(z-a)^n\right|
 \le M_\rho\frac{(r/\rho)^{N+1}}{1-r/\rho}.
\end{equation}
这是把从 \(N+1\) 开始的控制几何级数相加所得的，其中 \(f\) 表示原级数的和。先固定比值 \(r/\rho<1\)，再令 \(N\to\infty\)，误差便统一趋零；若让 \(r\uparrow\rho\)，分母变小，这项估计就不再提供统一的边界控制。

\ProseParagraphBreak

收敛半径还依赖展开中心。例如对 \(a\neq1\)，几何级数给出
\[
\frac1{1-z}=\frac1{1-a}\sum_{n=0}^\infty
\left(\frac{z-a}{1-a}\right)^n,
\quad |z-a|<|1-a|.
\]
由系数直接计算，半径恰为 \(|1-a|\)。同一个函数在不同中心的展开，能够覆盖的圆盘大小并不相同。

\subsection{逐项求导}
\label{b3:subsec:010202}

一致收敛只控制函数值，一般并不自动保留导数。例如实区间上的 \(g_N(t)=\sin(Nt)/N\) 一致趋于零，但 \(g_N'(0)=1\)，并不趋于极限函数在零点的导数。幂级数允许逐项求导，需要利用它在较小圆盘上的额外结构：微分所增加的 \(n,n(n-1),\ldots\) 能被几何衰减吸收。证明时先控制形式导数级数，再确认其和确实是原函数的导数。

\Needspace{6\baselineskip}
\begin{theorem}\label{b3:thm:01-power-differentiation}
设 \(a\in\mathbb C\)、\(c_n\in\mathbb C\ (n\ge0)\)，幂级数 \(f(z)=\sum_{n=0}^\infty c_n(z-a)^n\) 的收敛半径为 \(R\in(0,\infty]\)。则 \(f\) 在 \(D(a,R)\) 上全纯，并可任意次\term[zhuxiang-qiudao]{逐项求导}[termwise differentiation]：
\begin{equation}\label{b3:eq:01-power-derivatives}
f^{(k)}(z)=\sum_{n=k}^\infty\frac{n!}{(n-k)!}c_n(z-a)^{n-k}
\quad(k\in\mathbb Z,\ k\geq1).
\end{equation}
各导数级数在任意闭子圆盘上一致收敛，收敛半径仍为 \(R\)。特别地，\(c_k=f^{(k)}(a)/k!\)。
\end{theorem}
\begin{proof}
固定 \(0<r<\rho<R\)。由于级数在 \(a+\rho\) 收敛，通项有界，故存在 \(M\) 使 \(|c_n|\rho^n\leq M\)。于是对 \(|z-a|\leq r\)，有
\[
\left|\frac{n!}{(n-k)!}c_n(z-a)^{n-k}\right|
\leq M\rho^{-k}n^k\left(\frac r\rho\right)^{n-k}.
\]
右侧关于 \(n\geq k\) 可求和，这可由相邻项比值趋于 \(r/\rho<1\) 验证。因而所有形式导数级数在闭子盘上都一致收敛，其和连续。

先处理一阶。记部分和为 \(S_N\)，形式一阶导数级数的和为 \(g\)。当线段 \([z,z+h]\) 留在一个闭子盘中时，对多项式沿该线段积分得到
\[
S_N(z+h)-S_N(z)=h\int_0^1 S_N'(z+th)\,\mathrm{d}t.
\]
两组级数的一致收敛允许令 \(N\to\infty\)，故对 \(h\neq0\)，
\[
\frac{f(z+h)-f(z)}h=\int_0^1g(z+th)\,\mathrm{d}t.
\]
由 \(g\) 连续，右端在 \(h\to0\) 时趋于 \(g(z)\)。所以 \(f'=g\)。对导数级数重复同一论证，即得各阶公式。

估计已经说明导数级数的半径不小于 \(R\)。设一阶导数级数的半径为 \(R_1\)。对任意 \(0<s<R_1\)，该级数在距离中心为 \(s\) 的点绝对收敛，故
\[
\sum_{n=1}^\infty|c_n|s^n
\leq s\sum_{n=1}^\infty n|c_n|s^{n-1}<\infty.
\]
这说明每个 \(s<R_1\) 都不超过原级数的收敛半径，因而 \(R_1\le R\)。逐阶应用得到相同半径；当 \(R=\infty\) 时，前面的闭子盘估计已覆盖任意有限半径。最后代入 \(z=a\)，高于 \(k\) 次的项均为零，得到系数公式。
\end{proof}

公式\eqref{b3:eq:01-power-derivatives}也证明同一中心的幂级数表示唯一。若两个级数在中心的某个邻域内表示同一函数，对等式逐次求导并在中心求值，就得到每个系数相同。

\ProseParagraphBreak

例如在 \(|z|<1\) 内，对几何级数求导可得
\[
 \frac1{(1-z)^2}=\sum_{n=1}^{\infty}nz^{n-1}.
\]
两个级数的半径都为一，但这并不主张它们在边界有相同的收敛情况。相反，
\[
 L(z)=\sum_{n=1}^{\infty}\frac{z^n}{n^2}
\]
在闭单位圆盘上绝对一致收敛，而其形式导数在 \(z=1\) 变成调和级数。逐项求导定理适用于开圆盘中的每个点，并不把边界上的连续性提升为边界可微性。

\ProseParagraphBreak

还可以先逐项积分来构造一个导数已知的函数。若 \(f(z)=\sum_{n\ge0}c_n(z-a)^n\) 的半径为 \(R>0\)，则
\[
 F(z)=\sum_{n=0}^{\infty}\frac{c_n}{n+1}(z-a)^{n+1}
\]
在 \(D(a,R)\) 中收敛；对闭子盘，原级数的绝对控制再乘以有界距离即可控制新级数。刚证的求导定理给出 \(F'=f\)、\(F(a)=0\)。一般地，设 \(U\subseteq\mathbb C\) 为开集、\(h:U\to\mathbb C\)，满足 \(H'=h\) 的全纯函数 \(H:U\to\mathbb C\) 称为 \(h\) 的\term[yuanhanshu]{原函数}[primitive]。此处从已知幂级数构造了一个原函数，并没有断言一般开集上的全纯函数都有全局原函数。

\subsection{解析函数}
\label{b3:subsec:010203}

\Needspace{6\baselineskip}
\begin{definition}\label{b3:def:01-analytic}
设 \(\Omega\subseteq\mathbb C\) 为开集，\(f:\Omega\to\mathbb C\)。若对每个 \(a\in\Omega\)，都存在 \(r>0\)、\(D(a,r)\subseteq\Omega\) 及一个幂级数，使
\[
f(z)=\sum_{n=0}^\infty c_n(z-a)^n\quad(z\in D(a,r)),
\]
则称 \(f\) 在 \(\Omega\) 上\term[jiexi]{解析}[analytic]，也称它为\term[jiexi-hanshu]{解析函数}[analytic function]。
\end{definition}

解析性是每个点附近的局部性质，不要求在整个定义域上使用一个中心的级数。例如 \(1/(1-z)\) 在 \(\mathbb C\setminus\{1\}\) 的每个点都可用前面的几何级数展开，但以原点为中心的那一个级数只覆盖单位圆盘。定义要求的是可在各点重新展开，不是把一个展开式越过自己的收敛圆周。由\cref{b3:thm:01-power-differentiation}，解析函数必全纯，且它的局部展开就是以 \(f^{(n)}(a)/n!\) 为系数的\term[taylor-ji-shu]{Taylor 级数}[Taylor series]。

\Needspace{6\baselineskip}
\begin{proposition}\label{b3:prop:01-power-recenter}
设 \(a\in\mathbb C\)、\(c_n\in\mathbb C\)，幂级数 \(f(z)=\sum_{n=0}^{\infty}c_n(z-a)^n\) 的收敛半径为 \(R>0\)。则 \(f\) 在整个收敛圆盘 \(D(a,R)\) 上解析。更具体地，若 \(|b-a|<R\)，则对 \(|z-b|<R-|b-a|\)，
\[
f(z)=\sum_{k=0}^\infty\left(
\sum_{n=k}^\infty\binom nk c_n(b-a)^{n-k}\right)(z-b)^k.
\]
\end{proposition}
\begin{proof}
先固定 \(b\)，再取 \(0<r<R-|b-a|\)。对整个 \(|z-b|\le r\)，数 \(q=|b-a|+r\) 都小于 \(R\)。二项式公式给出同一个绝对值总和控制
\[
\sum_{n=0}^\infty\sum_{k=0}^n
\binom nk |c_n|\cdot|b-a|^{n-k}\cdot r^k
=\sum_{n=0}^\infty|c_n|q^n<\infty.
\]
每个固定 \(k\) 的系数级数也绝对收敛：上式中的第 \(k\) 列是可求和非负总和的一部分，再除以 \(r^k>0\) 即可。因此可以先按 \(n\) 展开 \(((b-a)+(z-b))^n\)，再按 \(k\) 合并：有限和重排不改变数值，绝对收敛使未计入部分的误差趋于零。误差控制对整个 \(|z-b|\le r\) 相同，故得到的展开在各闭子盘上一致收敛。内层系数按逐项微分公式等于 \(f^{(k)}(b)/k!\)，不依赖所选的 \(z\) 或 \(r\)。由于 \(r<R-|b-a|\) 任意，得到陈述中的开圆盘。这个半径只是由原收敛圆盘保证的下界；实际展开半径可以更大。
\end{proof}

实变量中的无限次可微并不保证局部 Taylor 展开等于原函数。考虑实函数
\[
 \psi(t)=
 \begin{cases}
 e^{-1/t^2},&t\ne0,\\
 0,&t=0.
 \end{cases}
\]
在 \(t\ne0\) 处，每阶导数都是 \(1/t\) 的某个多项式乘以 \(e^{-1/t^2}\)。对任意非负整数 \(m\)，有 \(|t|^{-m}e^{-1/t^2}\to0\)：令 \(s=1/t^2\to\infty\)，再由实指数级数中任意一个足够高次的正项控制 \(e^s\)，即可验证。依次把各阶导数在零点补为零，并用差商中多出的一个 \(1/t\) 因子作同样估计，便得到 \(\psi\in C^\infty(\mathbb R)\) 且 \(\psi^{(k)}(0)=0\) 对所有 \(k\ge0\) 成立。它的 Taylor 级数恒为零，函数却在 \(t\ne0\) 时严格为正。

\ProseParagraphBreak

复变量中要排除这种现象，必须把一个点的微分条件与它周围的函数值联系起来。下面的圆周积分公式承担这一任务；证明从复导数存在出发，不预设 \(f'\) 连续。

\subsection{全纯即解析}
\label{b3:subsec:010204}

三角形可以反复细分，因而能把一条边界的积分缩到局部微分余项所控制的尺度。圆盘的凸性又使两条线段之间恰好围出三角形，三角形积分消失便能给出圆盘内的原函数。接下来把这一方法用于差商，消去它在中心的一点例外，就得到圆周积分公式。\chapref{b3:ch:02}将把积分范围扩展到一般曲线，并讨论定义域的拓扑条件。

\ProseParagraphBreak

对分段 \(C^1\) 曲线 \(\gamma:[\alpha,\beta]\to\mathbb C\) 和其像上的连续函数 \(h\)，定义\term[fu-quxian-jifen]{复曲线积分}[complex line integral]为
\[
\int_\gamma h(\zeta)\,\mathrm{d}\zeta
\coloneqq\int_\alpha^\beta h\bigl(\gamma(t)\bigr)\cdot\gamma'(t)\,\mathrm{d}t.
\]
右端分为实部、虚部的一元积分。曲线长度为 \(\ell(\gamma)=\int_\alpha^\beta|\gamma'(t)|\,\mathrm{d}t\)，从而
\begin{equation}\label{b3:eq:01-contour-estimate}
\left|\int_\gamma h(\zeta)\,\mathrm{d}\zeta\right|
\leq\ell(\gamma)\cdot\max_{\zeta\in\gamma}|h(\zeta)|.
\end{equation}
反向行进使积分变号，连接曲线使积分相加；这些性质直接来自一元换元与积分可加性。本小节的三角形和圆周都取逆时针方向。
\symidx{\(\int_\gamma h(\zeta)\,\mathrm{d}\zeta\)：复曲线积分}

\ProseParagraphBreak

\term[goursat-dingli]{Goursat 定理}[Goursat's theorem]先给出三角形边界上的积分消失。

\Needspace{6\baselineskip}
\begin{lemma}[Goursat，三角形情形]\label{b3:lem:01-goursat-local}
设 \(T\subset\mathbb C\) 是一个非退化闭三角形，\(U\subset\mathbb C\) 为包含 \(T\) 的开集，\(f:U\to\mathbb C\) 全纯。边界 \(\partial T\) 取逆时针方向，则
\[
\int_{\partial T}f(\zeta)\,\mathrm{d}\zeta=0.
\]
\end{lemma}
\begin{proof}
将三边中点相连，把 \(T\) 分成四个相似小三角形。小三角形边界上的积分之和等于原积分，因为内部边以相反方向出现两次。故至少一个小三角形上的积分绝对值不小于原积分绝对值的四分之一。反复选择，得到嵌套闭三角形 \(T_n\)，其直径和周长均缩为原来的 \(2^{-n}\)，并有
\[
\left|\int_{\partial T}f\,\mathrm{d}\zeta\right|
\leq4^n\left|\int_{\partial T_n}f\,\mathrm{d}\zeta\right|.
\]
这些闭三角形非空且嵌套，并包含于紧集 \(T\)；紧性保证它们有公共点，直径趋零又保证公共点只有一个，记作 \(p\)。由于 \(U\) 包含整个闭三角形，即使 \(p\) 落在某条边上，\(f\) 仍在它的开邻域内复可微。在 \(p\) 处展开为
\[
f(\zeta)=f(p)+f'(p)(\zeta-p)+\varepsilon(\zeta)(\zeta-p),
\quad\varepsilon(\zeta)\to0\quad(\zeta\to p).
\]
在 \(p\) 处令 \(\varepsilon(p)=0\)。常数项与一次项在闭三角形上的积分为零，因为它们分别是 \(f(p)\zeta\) 和 \(f'(p)(\zeta-p)^2/2\) 的导数，沿各边使用一元微积分基本定理后端点项相消。对余项应用\eqref{b3:eq:01-contour-estimate}，得
\[
\left|\int_{\partial T}f\,\mathrm{d}\zeta\right|
\leq\ell(\partial T)\cdot\operatorname{diam}(T)
\cdot\sup_{\zeta\in T_n}|\varepsilon(\zeta)|.
\]
令 \(n\to\infty\)，右端趋于零，结论成立。
\end{proof}

细分使积分的可能损失为 \(4^n\)，而周长与直径各缩小 \(2^{-n}\)；三者恰好抵消，最后留下趋零的微分余项。这也解释了为什么只说“很小的三角形上积分很小”还不够：必须记录它相对于边长的阶数。证明只在所选极限点使用微分展开，没有对导数作积分，也未假定偏导数连续。

\Needspace{6\baselineskip}
\begin{lemma}\label{b3:lem:01-punctured-primitive}
设 \(D=D(a_0,R)\subset\mathbb C\)，其中 \(a_0\in\mathbb C\)、\(R>0\)，且 \(p\in D\)。若 \(h:D\to\mathbb C\) 连续，在 \(D\setminus\{p\}\) 上全纯，则存在全纯函数 \(H:D\to\mathbb C\) 满足 \(H'=h\)。因此 \(h\) 沿 \(D\) 内任意分段 \(C^1\) 闭曲线的积分为零。
\end{lemma}
\begin{proof}
先证明 \(h\) 沿每个闭三角形 \(T\subset D\) 的边界积分为零。把 \(T\) 均匀分成 \(4^n\) 个直径为原来 \(2^{-n}\) 的小三角形。不接触 \(p\) 的小三角形由\cref{b3:lem:01-goursat-local}积分为零。重复中点细分所得网格中，一个点若在小三角形内部，只属于一个闭三角形；若在边内，至多属于两个；若在顶点，至多属于六个。因而接触 \(p\) 的小三角形数目与细分次数无关。这些可能不为零的积分之和，绝对值不超过
\[
6\cdot2^{-n}\ell(\partial T)\cdot\max_T|h|.
\]
内部边仍两两抵消。\(h\) 连续保证最大值有限，令 \(n\to\infty\) 即得三角形积分为零。若 \(p\notin T\)，直接用前一引理即可。

固定 \(a\in D\)，试图把 \(h\) 沿路径积累起来。圆盘内每个点都能由一条线段从 \(a\) 到达，因此先定义 \(H(w)=\int_{[a,w]}h(\zeta)\,\mathrm{d}\zeta\)。圆盘凸性保证涉及的线段和三角形都留在 \(D\) 中。三角形积分为零给出
\[
H(w+v)-H(w)=\int_{[w,w+v]}h(\zeta)\,\mathrm{d}\zeta
=v\int_0^1h(w+tv)\,\mathrm{d}t.
\]
三点共线时，同一等式由线段积分的可加性得到。除以 \(v\neq0\)，由连续性令 \(v\to0\)，便有 \(H'(w)=h(w)\)。由于 \(H'=h\) 连续，\(H\) 作为实映射具有连续微分。沿每一段 \(C^1\) 曲线应用实参数链式法则，有 \((H\circ\gamma)'=h(\gamma)\gamma'\)；一元微积分基本定理使积分成为端点差，各段连接点的项相消，闭曲线的起点与终点又相同，所以积分为零。
\end{proof}

前一引理暂时允许一个点处未能确认复可微。它没有先声称 \(h\) 在该点也全纯，而是直接用连续性让接触该点的短边积分趋零；随后用三角形积分消失构造 \(H\)。这个余地使下一步能够处理差商在分母为零处的补定义。

\ProseParagraphBreak

接着建立\term[cauchy-jifen-gongshi]{Cauchy 积分公式}[Cauchy's integral formula]在圆盘内的形式。

\Needspace{6\baselineskip}
\begin{lemma}[Cauchy 积分公式，圆盘情形]\label{b3:lem:01-disk-cauchy}
设 \(a\in\mathbb C\)、\(R>0\)，\(f\) 在 \(D(a,R)\) 上全纯。若 \(0<r<R\)，则对 \(|z-a|<r\)，
\begin{equation}\label{b3:eq:01-disk-cauchy}
f(z)=\frac1{2\pi\mathrm{i}}\int_{|\zeta-a|=r}
\frac{f(\zeta)}{\zeta-z}\,\mathrm{d}\zeta.
\end{equation}
圆周取逆时针方向。
\end{lemma}
\begin{proof}
固定 \(z\)，令
\[
h(\zeta)=\begin{cases}
\dfrac{f(\zeta)-f(z)}{\zeta-z},&\zeta\neq z,\\[4pt]
f'(z),&\zeta=z.
\end{cases}
\]
当 \(\zeta\to z\) 时，商式按复导数定义趋于 \(f'(z)\)，所以 \(h\) 在这个特殊点连续；在其余点，连续性与全纯性都由商法则得到。由\cref{b3:lem:01-punctured-primitive}，\(h\) 的圆周积分为零。因此只需证明
\[
\int_{|\zeta-a|=r}\frac{\mathrm{d}\zeta}{\zeta-z}=2\pi\mathrm{i}.
\]
在该圆周上，几何级数
\[
\frac1{\zeta-z}=\sum_{n=0}^\infty
\frac{(z-a)^n}{(\zeta-a)^{n+1}}
\]
一致收敛，因为 \(|z-a|/r<1\)。具体地，前 \(N+1\) 项之后的尾项模在整个圆周上不超过
\[
 \frac1r\frac{(|z-a|/r)^{N+1}}{1-|z-a|/r},
\]
趋零速度与圆周上的点无关。由\eqref{b3:eq:01-contour-estimate}，积分误差至多再乘圆周长 \(2\pi r\)，故可逐项积分。\(n=0\) 项沿 \(\zeta=a+r(\cos t+\mathrm{i}\sin t)\)、\(0\leq t\leq2\pi\) 参数化，积分为 \(2\pi\mathrm{i}\)；\(n\geq1\) 时，\((\zeta-a)^{-n-1}\) 有原函数 \(-\frac1n(\zeta-a)^{-n}\)，其闭曲线积分为零。把结果代回 \(h\) 的零积分，得到公式。
\end{proof}

公式\eqref{b3:eq:01-disk-cauchy}表明，边界圆周上的函数值决定了圆盘内部的函数值。将核函数 \((\zeta-z)^{-1}\) 展开为几何级数，就能得到全纯函数的幂级数表示。

\Needspace{6\baselineskip}
\begin{theorem}[全纯与解析的等价性]\label{b3:thm:01-holomorphic-analytic}
设 \(U\subseteq\mathbb C\) 为开集，\(f:U\to\mathbb C\)。则 \(f\) 全纯，当且仅当它解析。若 \(a\in U\)、\(R>0\)、\(D(a,R)\subseteq U\)，且 \(f\) 全纯，则
\begin{equation}\label{b3:eq:01-holomorphic-taylor}
f(z)=\sum_{n=0}^\infty\frac{f^{(n)}(a)}{n!}(z-a)^n
\quad(|z-a|<R).
\end{equation}
此外，给定 \(0<\rho<r<R\)，令 \(M_r=\max_{|\zeta-a|=r}|f(\zeta)|\)。对每个整数 \(N\ge0\)，有
\begin{equation}\label{b3:eq:01-taylor-remainder}
 \sup_{|z-a|\le\rho}
 \left|f(z)-\sum_{n=0}^{N}\frac{f^{(n)}(a)}{n!}(z-a)^n\right|
 \le M_r\frac{(\rho/r)^{N+1}}{1-\rho/r}.
\end{equation}
\end{theorem}
\begin{proof}
解析推出全纯已由逐项求导证明。反之，设 \(f\) 全纯，固定 \(0<r<R\)。在\eqref{b3:eq:01-disk-cauchy}中展开刚才的几何级数；在任意 \(|z-a|\leq\rho<r\) 上，它关于圆周变量一致收敛，乘以连续有界的 \(f(\zeta)\) 后仍可逐项积分。于是
\[
f(z)=\sum_{n=0}^\infty c_n(z-a)^n,
\quad c_n=\frac1{2\pi\mathrm{i}}\int_{|\zeta-a|=r}
\frac{f(\zeta)}{(\zeta-a)^{n+1}}\,\mathrm{d}\zeta.
\]
若 \(M_r=\max_{|\zeta-a|=r}|f(\zeta)|\)，则 \(|c_n|\leq M_r/r^n\)，所以所得幂级数至少在半径 \(r\) 内收敛。由逐项求导及系数唯一性，\(c_n=f^{(n)}(a)/n!\)，不依赖选择的 \(r\)。任意 \(|z-a|<R\) 都可放在一个 \(r<R\) 的圆盘内，故得到\eqref{b3:eq:01-holomorphic-taylor}。由系数界，余项在 \(|z-a|\le\rho\) 上的模不超过 \(M_r\sum_{n>N}(\rho/r)^n\)，求和即得\eqref{b3:eq:01-taylor-remainder}。对一般开集，在每一点内取一个这样的圆盘即可。
\end{proof}

全纯函数因此拥有任意阶复导数，各阶导数仍全纯且连续。由 Cauchy--Riemann 关系，\(\partial_x f=f'\)、\(\partial_y f=\mathrm{i}f'\)；再对每个全纯的导数重复这个等式，可得
\[
 \partial_x^j\partial_y^k f=\mathrm{i}^{\,k}f^{(j+k)}
 \quad(j,k\ge0).
\]
所以实部、虚部也都是无限次连续可微的。幂级数表示把最初的一次复可微提升为任意阶连续可微。

\ProseParagraphBreak

在实际使用中，展开半径要区分“已保证的范围”和“最大的范围”。若只知道 \(f\) 在 \(D(a,R)\) 全纯，定理保证中心 Taylor 级数至少覆盖这个圆盘，不能推出半径恰为 \(R\)。例如同一函数 \(1/(1-z)\) 即使暂时只在 \(D(0,1/2)\) 上讨论，它的中心展开仍有半径一。另一方面，取 \(r<R\) 的圆周，是为了圆周本身处于定义域内部；不曾假定函数在 \(|z-a|=R\) 有连续边值。

\ProseParagraphBreak

本小节先证明了局部积分结论。\chapref{b3:ch:02}将在相应的拓扑条件下，把三角形上的结论推广到一般曲线，并系统讨论积分公式的推论；本章接下来则先用解析性构造初等函数、研究零点和局部映射的形状。

% !TeX root = ../../Book3.tex
\section{初等全纯函数}
\label{b3:sec:0103}

幂级数提供了构造全纯函数的直接方法。本节先把实变量的指数与三角函数延伸到复平面，再处理指数函数的局部逆。对数的困难来自不同点可以有同一个指数值：局部求逆没有障碍，沿闭曲线作全局选择却可能失败。最后讨论的 Möbius 变换则把代数运算与圆、直线的几何联系起来。

\subsection{指数函数}
\label{b3:subsec:010301}

\begin{definition}
\label{b3:def:01-elementary-exp}
对 \(z\in\mathbb C\)，令
\[
\exp z\coloneqq\sum_{n=0}^{\infty}\frac{z^n}{n!},
\quad e^z\coloneqq\exp z.
\]
这个函数称为复\term[zhishu-hanshu]{指数函数}[exponential function]。\symidx{\(\exp z=e^z\)：复指数函数}
\end{definition}

对每个固定的 \(R>0\)，相邻控制项 \(R^{n+1}/(n+1)!\) 与 \(R^n/n!\) 的比值趋于零。因此级数在每个闭圆盘上一致绝对收敛；由\cref{b3:thm:01-power-differentiation}，它在整个复平面上全纯，并可任意次逐项求导。实轴上的定义与实变量指数函数相同，这也可由实 Taylor 公式及其余项估计核对。下面的\term[euler-gongshi]{Euler 公式}[Euler's formula]进一步说明复指数与实三角函数之间的联系。

\begin{proposition}[Euler 公式]
\label{b3:prop:01-elementary-exp}
复指数函数满足
\[
(e^z)'=e^z,
\quad e^{z+w}=e^z\cdot e^w,
\quad e^z\neq0.
\]
若 \(x,y\in\mathbb R\)，则
\begin{equation}
\label{b3:eq:01-elementary-euler}
e^{x+\mathrm{i}y}=e^x\cdot(\cos y+\mathrm{i}\sin y),
\quad |e^{x+\mathrm{i}y}|=e^x.
\end{equation}
此外，\(e^z=e^w\) 当且仅当 \(z-w\in2\pi\mathrm{i}\mathbb Z\)。
\end{proposition}

\begin{proof}
逐项求导后把指标减一，得到导数公式。两个指数级数的绝对值级数均收敛，故其乘积可按总次数重新求和；总次数为 \(n\) 的项为
\[
\sum_{k=0}^{n}\frac{z^kw^{n-k}}{k!(n-k)!}
=\frac{(z+w)^n}{n!}.
\]
这里重新求和的合法性来自绝对尾项控制：保留足够大的有限矩形后，遗漏项绝对值之和可由两个收敛的非负级数的尾和控制到任意小。于是得到加法公式。令 \(w=-z\)，可知 \(e^z\cdot e^{-z}=1\)，从而指数函数没有零点。

把 \(e^{\mathrm{i}y}\) 的偶数次项与奇数次项分别收集，得到实余弦与实正弦的 Taylor 级数。实正弦、余弦的各阶导数绝对值不超过 \(1\)，Taylor 余项趋于零，因此
\[
e^{\mathrm{i}y}=\cos y+\mathrm{i}\sin y.
\]
再用加法公式及实恒等式 \(\cos^2y+\sin^2y=1\)，便得到\eqref{b3:eq:01-elementary-euler}。若 \(e^{z-w}=1\)，模长公式先给出 \(\operatorname{Re}(z-w)=0\)，实三角函数的取值再给出 \(\operatorname{Im}(z-w)\in2\pi\mathbb Z\)。反向结论直接代入即可。
\end{proof}

式\eqref{b3:eq:01-elementary-euler}把复指数分成径向伸缩与转动两部分。特别地，\(2\pi\mathrm{i}\) 是指数函数的周期，而指数函数的导数处处不为零。因此，导数不消失只能保证局部的可逆性，不能保证在整个定义域上单射。

\subsection{对数函数}
\label{b3:subsec:010302}

对非零复数 \(z\)，方程 \(e^w=z\) 的解的实部必须等于 \(\log|z|\)，虚部则是 \(z\) 的一个辐角；所有解恰好相差 \(2\pi\mathrm{i}\) 的整数倍。要把这些解组成一个函数，需要连续地选取辐角。

\begin{definition}
\label{b3:def:01-elementary-log}
在割平面 \(D\coloneqq\mathbb C\setminus(-\infty,0]\) 上，取满足
\[
z=|z|e^{\mathrm{i}\theta},\quad -\pi<\theta<\pi
\]
的唯一实数 \(\theta\)，记为 \(\operatorname{Arg}z\)，称为\term[fujiao-zhuzhi]{辐角主值}[principal argument]。函数
\[
\operatorname{Log}z\coloneqq\log|z|+\mathrm{i}\operatorname{Arg}z
\]
称为\term[duishu-zhuzhi]{对数主值}[principal logarithm]。\symidx{\(\operatorname{Arg}z\)：取值于 \((-\pi,\pi)\) 的辐角主值}\symidx{\(\operatorname{Log}z\)：割平面上的对数主值}
\end{definition}

这里实数的对数仍记为 \(\log\)。本书把负实轴连同原点从主值的定义域中删去；另有把非零复数的辐角主值取在 \((-\pi,\pi]\) 的约定。\footnote{参见王以忠《工程复变函数与积分变换》（清华大学出版社，2021）\href{https://www.tup.tsinghua.edu.cn/upload/books/yz/091462-01.pdf}{出版社样章}印刷第 8 页。该处使用 \((-\pi,\pi]\) 选取点值；本节为讨论全纯分支而删去负实轴。}给负实数指定端点辐角，不能使主值在那里连续。

\begin{proposition}
\label{b3:prop:01-elementary-principal-log}
对数主值在 \(D\) 上全纯，并且
\[
\operatorname{Log}1=0,
\quad e^{\operatorname{Log}z}=z,
\quad (\operatorname{Log}z)'=\frac1z.
\]
当且仅当 \(-\pi<\operatorname{Im}w<\pi\) 时，有 \(e^w\in D\) 且 \(\operatorname{Log}(e^w)=w\)。
\end{proposition}

\begin{proof}
写 \(z=x+\mathrm{i}y\)、\(r=(x^2+y^2)^{1/2}\)。在 \(D\) 上有 \(r+x>0\)，实半角公式给出
\[
\operatorname{Arg}(x+\mathrm{i}y)
=2\arctan\frac{y}{r+x}.
\]
右边取值于 \((-\pi,\pi)\)，其余弦为 \(x/r\)、正弦为 \(y/r\)，所以确实是定义中的唯一辐角。这个表达式还证明辐角主值是光滑的。对其求实偏导，得到
\[
(\operatorname{Arg}z)_x=-\frac{y}{r^2},
\quad (\operatorname{Arg}z)_y=\frac{x}{r^2}.
\]
另一方面，\(\log r\) 的两个偏导为 \(x/r^2\) 与 \(y/r^2\)。它们满足 Cauchy--Riemann 方程，由\cref{b3:cor:01-continuous-partial-criterion}可知主对数全纯，且其复导数为
\[
\frac{x}{r^2}-\mathrm{i}\frac{y}{r^2}=\frac1z.
\]
其余断言由定义及\cref{b3:prop:01-elementary-exp}得出：指数值唯一确定实部，并把虚部确定到模 \(2\pi\)；主值恰选取位于 \((-\pi,\pi)\) 的代表。
\end{proof}

主值不能无条件保留实对数的乘法规则。例如取 \(z=e^{3\pi\mathrm{i}/4}\)，则 \(z^2=-\mathrm{i}\)，所以
\[
\operatorname{Log}(z^2)=-\frac{\pi\mathrm{i}}2,
\quad 2\operatorname{Log}z=\frac{3\pi\mathrm{i}}2.
\]
两者相差一个完整的周期。类似地，当正数 \(r\) 固定时，从负实轴的上侧和下侧趋近 \(-r\)，主对数分别趋于 \(\log r+\pi\mathrm{i}\) 与 \(\log r-\pi\mathrm{i}\)。这解释了为什么不能把两侧的主值直接接起来。

\begin{definition}
\label{b3:def:01-elementary-log-branch}
设 \(U\subseteq\mathbb C\setminus\{0\}\) 是开集。若连续函数 \(L\colon U\rightarrow\mathbb C\) 满足 \(e^{L(z)}=z\)，则称 \(L\) 为 \(U\) 上的一个\term[duishu-fenzhi]{对数分支}[branch of the logarithm]。
\end{definition}

\begin{proposition}
\label{b3:prop:01-elementary-local-log}
每个非零点都有一个邻域，在该邻域上存在全纯的对数分支。任何连续的对数分支自动全纯，导数为 \(1/z\)；在连通开集上，两个分支之差是某个固定的 \(2\pi\mathrm{i}k\)，其中 \(k\in\mathbb Z\)。但整个 \(\mathbb C\setminus\{0\}\) 上不存在连续的对数分支。
\end{proposition}

\begin{proof}
固定 \(z_0\neq0\)，选取 \(\ell_0\) 使 \(e^{\ell_0}=z_0\)。当 \(|z-z_0|<|z_0|\) 时，\(z/z_0\) 的实部为正，故
\[
L_0(z)=\ell_0+\operatorname{Log}(z/z_0)
\]
有定义，且是所需的局部分支。若 \(L\) 是同一邻域上的任意连续分支，则 \((L-L_0)/(2\pi\mathrm{i})\) 是整数值连续函数，在较小的连通邻域上恒定。因此 \(L\) 局部全纯且 \(L'=1/z\)。同一理由说明连通开集上两个分支的差恒定。

若整个去心平面上存在连续分支 \(L\)，令
\[
h(t)=\frac{L(e^{\mathrm{i}t})-\mathrm{i}t}{2\pi\mathrm{i}},
\quad 0\leq t\leq2\pi.
\]
由指数函数的周期刻画，\(h\) 是连续的整数值函数，因而恒定。但 \(e^{2\pi\mathrm{i}}=1\) 使 \(h(2\pi)=h(0)-1\)，矛盾。
\end{proof}

有了一个分支 \(L\)，便可在其定义域上令 \(z^{\alpha}=\exp(\alpha L(z))\)。改变分支会把这个值乘以 \(e^{2\pi\mathrm{i}k\alpha}\)。整数幂与分支无关，平方根等非整数幂则必须交代所选分支。\subsecref{b3:subsec:010404}构造一个非零全纯函数的局部根时，只需在像靠近某个非零值的小邻域内使用这里的局部分支。

\subsection{三角函数}
\label{b3:subsec:010303}

\begin{definition}
\label{b3:def:01-elementary-trigonometric}
在复平面上定义\term[zhengxian-hanshu]{正弦函数}[sine function]与\term[yuxian-hanshu]{余弦函数}[cosine function]为
\[
\sin z\coloneqq\frac{e^{\mathrm{i}z}-e^{-\mathrm{i}z}}{2\mathrm{i}},
\quad \cos z\coloneqq\frac{e^{\mathrm{i}z}+e^{-\mathrm{i}z}}2.
\]
\symidx{\(\sin z\)：复正弦函数}\symidx{\(\cos z\)：复余弦函数}
\end{definition}

由 Euler 公式，这些定义在实轴上还原通常的正弦与余弦。求导及相乘给出
\[
(\sin z)'=\cos z,
\quad (\cos z)'=-\sin z,
\quad \sin^2z+\cos^2z=1.
\]
例如最后一个等式来自
\[
\frac{(e^{\mathrm{i}z}+e^{-\mathrm{i}z})^2
-(e^{\mathrm{i}z}-e^{-\mathrm{i}z})^2}{4}=1.
\]
把指数的加法公式代入定义，还得到
\[
\begin{aligned}
\sin(z+w)&=\sin z\cdot\cos w+\cos z\cdot\sin w,\\
\cos(z+w)&=\cos z\cdot\cos w-\sin z\cdot\sin w.
\end{aligned}
\]
同样，由指数级数的奇偶项分别得到
\[
\sin z=\sum_{n=0}^{\infty}\frac{(-1)^nz^{2n+1}}{(2n+1)!},
\quad
\cos z=\sum_{n=0}^{\infty}\frac{(-1)^nz^{2n}}{(2n)!}.
\]
它们都是整个复平面上的全纯函数。由指数函数的周期刻画，\(\sin z=0\) 恰在 \(z\in\pi\mathbb Z\) 时发生，\(\cos z=0\) 恰在 \(z\in\pi/2+\pi\mathbb Z\) 时发生；例如前者等价于 \(e^{2\mathrm{i}z}=1\)，后者等价于 \(e^{2\mathrm{i}z}=-1\)。

\ProseParagraphBreak
复三角函数保留上述恒等式，却不保留实轴上的有界性。沿虚轴有
\[
\sin(\mathrm{i}y)=\frac{\mathrm{i}}2(e^y-e^{-y}),
\quad \cos(\mathrm{i}y)=\frac12(e^y+e^{-y}),
\quad y\in\mathbb R.
\]
当 \(y\rightarrow+\infty\) 时，两者的模都以指数速度增长。恒等式 \(\sin^2z+\cos^2z=1\) 中的平方没有复共轭，因而不能推出 \(|\sin z|^2+|\cos z|^2=1\)。计算复函数的大小时，代数恒等式与含模估计必须分别处理。

\subsection{Möbius 变换}
\label{b3:subsec:010304}

分式线性函数可能在一个有限点没有有限值。把这个值解释为无穷远点以后，它们可以作为整个扩充平面上的可逆映射来讨论。

\begin{definition}
\label{b3:def:01-elementary-extended-plane}
在复平面上添加一个点 \(\infty\)，得到\term[kuochong-fu-pingmian]{扩充复平面}[extended complex plane]
\[
% 零宽核避免 XeTeX 对 AMS 单字母加 unicode-math 重音时的字体错误。
\widehat{\mathbb C\kern0pt}\coloneqq\mathbb C\cup\{\infty\}.
\]
\symidx{\(\widehat{\mathbb C\kern0pt}\)：扩充复平面}
无穷远点的基本邻域取为 \(\{\infty\}\cup\{z:|z|>R\}\)，其中 \(R>0\)。在这样的邻域内使用坐标 \(\zeta=1/z\)，并规定无穷远点对应 \(\zeta=0\)。
\end{definition}

讨论映射在无穷远点附近的连续性或全纯性，是指改用上述坐标以后讨论相应的普通函数；若函数值也为无穷远，则目标同样使用倒数坐标。例如 \(z\mapsto1/z\) 在扩充平面上交换 \(0\) 与 \(\infty\)，在相应坐标中就是恒等函数。本节只需这两张坐标图，不预先使用一般曲面的理论。

\begin{definition}
\label{b3:def:01-elementary-mobius}
设 \(a,b,c,d\in\mathbb C\) 且 \(ad-bc\neq0\)。在分母不为零处定义
\[
T(z)=\frac{az+b}{cz+d}.
\]
若 \(c\neq0\)，补定义 \(T(-d/c)=\infty\)、\(T(\infty)=a/c\)；若 \(c=0\)，补定义 \(T(\infty)=\infty\)。这样得到的扩充平面上的映射称为\term[mobius-bianhuan]{Möbius 变换}[Möbius transformation]，也称分式线性变换。\footnote{“分式线性变换”见王晓光《复变函数》（高等教育出版社，2025）第十三章的\href{https://xuanshu.hep.com.cn/front/h5Mobile/bookDetails?bookId=68598849e119ac9729fe438b}{出版社目录}。}\termsee[fenshi-xianxing-bianhuan]{分式线性变换}{Möbius 变换}
\end{definition}

\begin{proposition}
\label{b3:prop:01-elementary-mobius-inverse}
每个 Möbius 变换都是扩充平面上的连续双射，其逆变换为
\[
T^{-1}(w)=\frac{dw-b}{-cw+a},
\]
其中无穷值按定义中的规则解释。在有限且函数值有限的点处，
\[
T'(z)=\frac{ad-bc}{(cz+d)^2}\neq0.
\]
在源与目标的其他坐标中，它也全纯且导数不为零。
\end{proposition}

\begin{proof}
令 \(\Delta=ad-bc\)。从 \(w(cz+d)=az+b\) 解出 \(z\)，得到逆变换公式；其系数行列式也是 \(\Delta\neq0\)，且代入补定义的特殊点仍互为逆。因此映射是双射，普通点处的连续性和导数公式由商法则得到。

若 \(c\neq0\)，在 \(z_0=-d/c\) 附近给目标使用倒数坐标，所得函数为
\[
\frac1{T(z)}=\frac{cz+d}{az+b}.
\]
分母在 \(z_0\) 处等于 \(-\Delta/c\neq0\)，导数在该点为 \(-c^2/\Delta\neq0\)。在源点 \(\infty\) 附近令 \(z=1/\zeta\)，则
\[
T(1/\zeta)=\frac{a+b\zeta}{c+d\zeta},
\]
右边在 \(\zeta=0\) 处全纯，导数为 \(-\Delta/c^2\neq0\)。若 \(c=0\)，则 \(a,d\neq0\)，在源与目标同时使用倒数坐标，得到
\[
\frac1{T(1/\zeta)}=\frac{d\zeta}{a+b\zeta},
\]
其在零点处的导数为 \(d/a\neq0\)。这些式子同时证明所补无穷值处的连续性及局部全纯性。
\end{proof}

下面把圆与直线放在一起讨论。圆取通常的有限半径圆，直线则补入无穷远点；这一约定使 Möbius 变换的几何性质能够统一表述。

\begin{proposition}
\label{b3:prop:01-elementary-mobius-circles}
Möbius 变换把每个圆或直线双射到一个圆或直线。
\end{proposition}

\begin{proof}
圆或直线的有限部分统一写成
\[
A|z|^2+Bz+\overline B\,\overline z+C=0,
\quad A,C\in\mathbb R,
\quad |B|^2-AC>0.
\]
若 \(A\neq0\)，配方后这是圆心为 \(-\overline B/A\)、半径为 \(\sqrt{|B|^2-AC}/|A|\) 的圆；若 \(A=0\)，则 \(B\neq0\)，方程给出一条直线，并按约定补入 \(\infty\)。

平移和非零复数乘法把圆变为圆、直线变为直线：例如代入 \(z=(w-\beta)/\alpha\)、乘以 \(|\alpha|^2\)，所得方程仍具同一形式，且新系数的判别量为 \(|\alpha|^2\cdot\bigl(|B|^2-AC\bigr)>0\)。对于倒数变换 \(w=1/z\)，在 \(w\neq0\) 处代入并乘以 \(|w|^2\)，得到
\[
C|w|^2+\overline B w+B\overline w+A=0.
\]
判别量没有改变。原方程含有 \(0\) 当且仅当 \(C=0\)，恰与像方程表示直线、因而含有 \(\infty\) 对应；原方程表示直线当且仅当 \(A=0\)，恰与像中含有 \(0\) 对应。所以特殊点也包含在这个对应中。

当 \(c=0\) 时，\(T\) 本身是非退化的仿射变换；当 \(c\neq0\) 时，可以分解为
\[
T(z)=\frac ac-\frac{ad-bc}{c^2}\cdot\frac1{z+d/c}.
\]
这是平移、倒数变换、非零复数乘法与平移的复合。前面的计算证明其像是圆或直线，逆变换又保证映射到的是整个圆或整条直线。
\end{proof}

一个常用例子是
\[
T(z)=\frac{z-\mathrm{i}}{z+\mathrm{i}}.
\]
当 \(\operatorname{Im}z>0\) 时，\(|z+\mathrm{i}|^2-|z-\mathrm{i}|^2=4\operatorname{Im}z>0\)，故 \(|T(z)|<1\)。反过来，由 \(z=\mathrm{i}(1+w)/(1-w)\) 算出
\[
\operatorname{Im}z=\frac{1-|w|^2}{|1-w|^2}>0
\quad (|w|<1).
\]
因此这个变换把上半平面双射到单位圆盘；它把扩充实轴送到单位圆周，无穷远点对应圆周上的 \(1\)。借助这个映射，可以将半平面上的问题化为圆盘上的问题，或反过来处理。

% !TeX root = ../../Book3.tex
\section{局部结构}
\label{b3:sec:0104}

局部幂级数把函数在一点附近的行为分成两部分：最先出现的非零项确定主要形状，其后的项给出可控制的修正。先用这种分解研究零点，再把局部相等沿连通区域传开；最后选择新的局部坐标，将非常值全纯函数化为一个幂映射。这样既能解释导数不为零时的可逆性，也能看清导数为零时原像怎样合并。

\subsection{孤立零点}
\label{b3:subsec:010401}

\Needspace{6\baselineskip}
\begin{definition}\label{b3:def:01-zero-order}
设 \(\Omega\subseteq\mathbb C\) 为开集，\(a\in\Omega\)，\(f:\Omega\to\mathbb C\) 全纯且 \(f(a)=0\)。若存在正整数 \(m\)，使
\[
f(a)=f'(a)=\cdots=f^{(m-1)}(a)=0,
\quad f^{(m)}(a)\neq0,
\]
则称 \(a\) 是 \(f\) 的 \(m\) 阶零点，\(m\) 称为\term[lingdian-jieshu]{零点阶数}[order of a zero]，记为 \(\operatorname{ord}_a f\)。\(m=1\) 时称为\term[dan-lingdian]{单零点}[simple zero]。若 \(a\) 的某邻域内没有 \(f\) 的其他零点，则称 \(a\) 为\term[guli-lingdian]{孤立零点}[isolated zero]。
\end{definition}
\symidx{\(\operatorname{ord}_a f\)：全纯函数在 \(a\) 处的零点阶数}

\Needspace{6\baselineskip}
\begin{proposition}\label{b3:prop:01-zero-factorization}
设 \(\Omega\subseteq\mathbb C\) 为开集，\(a\in\Omega\)，\(f:\Omega\to\mathbb C\) 全纯且 \(f(a)=0\)。则恰有以下两种情形之一：\(f\) 在 \(a\) 的一个邻域内恒为零；或存在唯一正整数 \(m\)、圆盘 \(D(a,R)\subseteq\Omega\) 及其上的全纯函数 \(g\)，使
\begin{equation}\label{b3:eq:01-zero-factorization}
f(z)=(z-a)^m\cdot g(z),
\quad g(a)\neq0.
\end{equation}
后一种情形中，\(a\) 是孤立零点，且阶数为 \(m\)；在所取圆盘上，函数 \(g\) 也由 \(f\) 唯一确定。
\end{proposition}
\begin{proof}
取以 \(a\) 为中心、包含于 \(\Omega\) 的圆盘，在其中使用\eqref{b3:eq:01-holomorphic-taylor}。若所有系数都为零，函数就在该圆盘内恒为零。否则，非零系数的指标组成非空的非负整数集，有最小元；又因 \(f(a)=0\)，这个最小元 \(m\) 至少为一。提出 \((z-a)^m\)，余下级数
\[
g(z)=\sum_{k=0}^\infty\frac{f^{(m+k)}(a)}{(m+k)!}(z-a)^k
\]
在原圆盘内部的每个较小闭圆盘上绝对一致收敛，因而全纯，并且 \(g(a)=f^{(m)}(a)/m!\neq0\)。由连续性，在更小的圆盘内有 \(|g(z)|\geq|g(a)|/2>0\)，所以其中只有 \(a\) 一个零点。系数唯一性确定 \(m\)；在 \(z\neq a\) 处 \(g\) 由商式唯一确定，在 \(a\) 处再由连续性确定。
\end{proof}

例如 \(e^z-1\) 在原点的一阶系数为 \(1\)，所以原点是单零点；\(1-\cos z=z^2/2+\cdots\) 在原点的零点阶数为 \(2\)。若两个全纯函数 \(f,g\) 在 \(a\) 的零点阶数分别为 \(m,n\)，相乘其局部分解便有
\[
\operatorname{ord}_a(fg)=m+n.
\]
求和则可能发生首项抵消。例如 \(e^z-1\) 和 \(-z\) 都在原点一阶消失，它们的和 \(e^z-1-z\) 却在原点二阶消失。阶数记录的是实际留下的第一个非零项，不能只根据两个输入阶数决定和的阶数。

\ProseParagraphBreak

阶数还给出邻域内的双边估计。若 \(f(z)=(z-a)^m\cdot g(z)\)、\(g(a)\ne0\)，由连续性缩小半径，可使
\[
 \frac{|g(a)|}{2}\cdot|z-a|^m\le |f(z)|
 \le\frac{3|g(a)|}{2}\cdot|z-a|^m.
\]
常数对这个圆盘中的所有点同时有效，所以有限阶零点沿任意方向都以同一整数次幂消失。以 \(e^z-1-z\) 为例，除以 \(z^2\) 后的全纯因子在原点取值 \(1/2\)，因而其模在充分小的复圆盘内与 \(|z|^2\) 双边可比。

\ProseParagraphBreak

反过来，设 \(f\) 在 \(a\) 附近全纯，\(m\) 为非负整数，\(\alpha>0\)。若存在常数 \(C>0\)，使在趋于 \(a\) 的一列不同点 \(z_n\ne a\) 上有
\[
 |f(z_n)|\le C\cdot|z_n-a|^{m+\alpha},
\]
则连续性先给 \(f(a)=0\)。若它的有限零点阶数为 \(k\le m\)，写 \(f(z)=(z-a)^k\cdot g(z)\)、\(g(a)\ne0\)，上面的下界会给出
\[
 0<\frac{|g(a)|}{2}
 \le C\cdot|z_n-a|^{m+\alpha-k}\longrightarrow0,
\]
矛盾。因此 \(f\) 的阶数至少为 \(m+1\)，或它局部恒零。特别地，对已经知道全纯的函数，沿实轴趋近取得这样的估计也足够；这里的依据正是全纯性事先给出的邻域分解。对一般实可微函数则不能这样推断，例如 \(z\mapsto\operatorname{Im}z\) 在实轴恒为零，在虚方向却只一阶消失。单条路径上的信息与整个邻域的统一结论之间，必须有相应的结构保证。

\subsection{恒等定理}
\label{b3:subsec:010402}

一个零点的附近可以由幂级数控制；要把局部信息传遍定义域，还需要连通性。\term[hengdeng-dingli]{恒等定理}[identity theorem]，也称唯一性定理或惟一性定理，\footnote{王晓光《复变函数》（高等教育出版社，2025）第 11.3 节的\href{https://xuanshu.hep.com.cn/front/h5Mobile/bookDetails?bookId=68598849e119ac9729fe438b}{出版社目录}用“唯一性定理”；钟玉泉《复变函数论》（第五版，高等教育出版社，2021）第四章第 4 节的\href{https://xuanshu.hep.com.cn/front/book/findBookDetails?bookId=60202746966feeeca73340f0}{出版社目录}用“惟一性定理”。}\termsee[weiyixing-dingli]{唯一性定理}{恒等定理}\termsee[weiyixing-dingli-old]{惟一性定理}{恒等定理}说明，在区域内部积累的一列相等函数值，已经足以确定整个全纯函数。

\Needspace{6\baselineskip}
\begin{theorem}[恒等定理]\label{b3:thm:01-identity}
设 \(\Omega\subseteq\mathbb C\) 为区域，\(f,g:\Omega\to\mathbb C\) 全纯。若集合 \(\{z\in\Omega:f(z)=g(z)\}\) 在 \(\Omega\) 内有一个聚点，则 \(f=g\) 于整个 \(\Omega\)。
\end{theorem}
\begin{proof}
令 \(h=f-g\)，设所述聚点为 \(a\in\Omega\)。由连续性，\(h(a)=0\)。\Cref{b3:prop:01-zero-factorization}排除了有限阶零点，因为那会使 \(a\) 附近没有其他零点。因此 \(h\) 在 \(a\) 的某邻域内恒为零。

令 \(E\) 为 \(\Omega\) 中那些具有恒零邻域的点组成的集合。它非空且开；由 Taylor 展开，亦可写为
\[
E=\bigcap_{k=0}^\infty\{z\in\Omega:h^{(k)}(z)=0\}.
\]
确实，恒零邻域使所有导数消失，反过来所有 Taylor 系数为零又给出恒零邻域。各阶导数连续，所以 \(E\) 在 \(\Omega\) 中闭。区域连通性使非空的开闭子集 \(E\) 等于 \(\Omega\)，结论得证。
\end{proof}

\Needspace{6\baselineskip}
\begin{corollary}\label{b3:cor:01-discrete-zeros}
设 \(\Omega\subseteq\mathbb C\) 为区域，\(f:\Omega\to\mathbb C\) 全纯且不恒为零。则 \(f\) 的每个零点都孤立，并且在每个紧集 \(K\subseteq\Omega\) 内只有有限个零点。
\end{corollary}
\begin{proof}
若有零点不是有限阶，\cref{b3:prop:01-zero-factorization}给出恒零邻域，再由恒等定理得到全域恒零，矛盾。所以每个零点都孤立。若紧集内有无穷多个不同零点，紧致性给出其中一列的聚点；聚点仍在 \(K\subseteq\Omega\) 内，再次与恒等定理矛盾。
\end{proof}

聚点位于区域内部这一条件不能删去。函数 \(\sin(\pi/z)\) 在 \(\mathbb C\setminus\{0\}\) 上全纯，具有零点 \(1/n\)、\(n\geq1\)，却不恒为零；这些零点的聚点 \(0\) 不在定义域内。连通性也有自己的作用：在两个不交圆盘上分别取常数 \(0\) 和 \(1\)，得到一个全纯函数，它在一个圆盘内恒为零，却不能因此确定另一个圆盘上的值。

\ProseParagraphBreak

一个常用后果是：若两个全纯函数在区域内的一段非退化实区间上相等，则它们在整个区域相等。实区间在复平面中没有内部，但它拥有位于定义域内的聚点。实际判断时甚至不需要整段区间：若整函数 \(f\) 满足 \(f(1/n)=e^{1/n}\) 对所有正整数 \(n\) 成立，聚点 \(0\) 位于共同定义域内，便有 \(f(z)=e^z\) 于全平面。相比之下，只给 \(f(n)=e^n\) 并不能唯一确定 \(f\)，因为 \(e^z+\sin\pi z\) 也满足同样条件，而正整数在平面内没有聚点。

\subsection{局部逆函数定理}
\label{b3:subsec:010403}

复导数不为零时，微分是可逆的旋转伸缩。局部\term[ni-hanshu-dingli]{逆函数定理}[inverse function theorem]把这一点上的可逆性提升为邻域内函数的可逆性。本书沿用第一、二卷的“逆函数”主称。

\Needspace{6\baselineskip}
\begin{theorem}[局部逆函数]\label{b3:thm:01-local-inverse}
设 \(\Omega\subseteq\mathbb C\) 为开集，\(a\in\Omega\)，\(f:\Omega\to\mathbb C\) 全纯且 \(f'(a)\neq0\)。则存在开邻域 \(U\subseteq\Omega\)、\(V\subseteq\mathbb C\)，分别包含 \(a\)、\(f(a)\)，使 \(f:U\to V\) 为双射，\(f'\) 在 \(U\) 上不为零，其逆函数 \(h:V\to U\) 全纯，并满足
\[
h'(w)=\frac1{f'\bigl(h(w)\bigr)}\quad(w\in V).
\]
\end{theorem}
\begin{proof}
由\cref{b3:thm:01-holomorphic-analytic}，将 \(f\) 看作 \(\mathbb R^2\) 间的映射时，它为 \(C^1\)。其在 \(a\) 处的实微分为乘以 \(f'(a)\) 的复线性映射，Jacobi 行列式为 \(|f'(a)|^2>0\)。第一卷引理 18.1.2 给出的实逆函数定理于是提供邻域间的 \(C^1\) 逆映射。由 \(f'\) 的连续性，可在所得源邻域内取更小的开邻域 \(U\)，使 \(f'\) 处处非零。实逆映射连续，故 \(V=f(U)\) 是原目标邻域中的开集；限制映射后仍得到邻域间的双射及连续逆映射 \(h\)。

固定 \(w_0\in V\)，记 \(z_0=h(w_0)\)。当 \(w\neq w_0\) 且 \(w\to w_0\) 时，单射性与逆映射连续性分别给出 \(z=h(w)\ne z_0\) 及 \(z\to z_0\)，因此
\[
\frac{h(w)-h(w_0)}{w-w_0}
=\left(\frac{f(z)-f(z_0)}{z-z_0}\right)^{-1}
\longrightarrow\frac1{f'(z_0)}.
\]
所以逆映射在每个点复可微，得到结论。
\end{proof}
% 跨卷引用：Book1/Part03/Chapter18/Section1801.tex，b1:lem:level-set-local-inverse，当前引理18.1.2；完整实逆函数证明不在本卷重复。

例如 \(f(z)=z^2\) 在 \(a=1\) 的导数为 \(2\)。目标点 \(w=1\) 附近选择
\[
 h(w)=\exp\bigl(\tfrac12\operatorname{Log}w\bigr),
\]
便有 \(h(1)=1\)、\(h(w)^2=w\) 和 \(h'(1)=1/2\)。另一个解 \(-h(w)\) 靠近 \(-1\)，不属于这里选定的源邻域。所谓局部逆函数，既要指定目标点附近的范围，也要指定逆值靠近哪个原像。

\ProseParagraphBreak

指数函数处处导数非零，因此每一点附近都可求全纯逆函数；但 \(e^{z+2\pi\mathrm{i}}=e^z\)，它在整个复平面上并非单射。若固定 \(e^{z_0}=w_0\ne0\)，则
\[
 h(w)=z_0+\operatorname{Log}(w/w_0)
 \quad\bigl(|w-w_0|<|w_0|\bigr)
\]
给出满足 \(h(w_0)=z_0\) 的逆分支。这里 \(w/w_0\) 落在右半平面，主对数有定义，且 \(e^{h(w)}=w\)。把 \(z_0\) 换成 \(z_0+2\pi\mathrm{i}k\) 就得到其他分支；局部求逆不会消除全局的周期。

\subsection{开映射定理}
\label{b3:subsec:010404}

在导数为零的点，局部逆函数定理不再适用。零点分解仍给出
\(f(z)-f(a)=(z-a)^m\cdot g(z)\)，其中 \(g\) 在小邻域内无零。若能把这个非零因子写成 \(q^m\)，便可令 \(\phi=(z-a)q\)，把所有高阶修正吸收到一个新坐标中。需要求根的是接近非零常数的 \(g\)，因此前节的局部对数已经足够。

\Needspace{6\baselineskip}
\begin{proposition}[局部标准形]\label{b3:prop:01-local-normal-form}
设 \(\Omega\subseteq\mathbb C\) 为开集，\(a\in\Omega\)，\(f:\Omega\to\mathbb C\) 全纯，且 \(f-f(a)\) 不在 \(a\) 的任何邻域内恒为零。令 \(m\geq1\) 为 \(f-f(a)\) 在 \(a\) 的零点阶数。则存在开邻域 \(U\subseteq\Omega\)、\(a\in U\)，数 \(r>0\) 及全纯双射 \(\phi:U\to D(0,r)\)，其逆函数也全纯，且 \(\phi(a)=0\)、\(\phi'\) 在 \(U\) 上不为零，并有
\begin{equation}\label{b3:eq:01-local-normal-form}
f(z)=f(a)+\phi(z)^m\quad(z\in U).
\end{equation}
邻域 \(U\) 可以取得包含于任意预先给定的 \(a\) 的开邻域。
\end{proposition}
\begin{proof}
由零点分解，在一个足够小的圆盘 \(D(a,R)\subseteq\Omega\) 上有 \(f(z)-f(a)=(z-a)^m\cdot g(z)\)，其中 \(g\) 全纯且 \(g(a)\neq0\)。若预先给定了 \(a\) 的开邻域，开始时便把这个圆盘取在其中。取 \(c\in\mathbb C\) 满足 \(c^m=g(a)\)；由非零复数的极坐标表示可取得这样的 \(c\)，且 \(c\ne0\)。继续缩小 \(R\)，使
\[
 \left|\frac{g(z)}{g(a)}-1\right|<\frac12
 \quad\bigl(z\in D(a,R)\bigr).
\]
因此 \(g(z)/g(a)\) 始终落在主对数的定义域内。利用\cref{b3:prop:01-elementary-principal-log}，在这个圆盘上定义
\[
 q(z)=c\cdot\exp\left(\frac1m\operatorname{Log}\frac{g(z)}{g(a)}\right).
\]
指数和对数的恒等式给出 \(q(a)=c\)、\(q(z)^m=g(z)\)。于是全纯函数 \(\phi(z)=(z-a)\cdot q(z)\) 满足 \(\phi(a)=0\)、\(\phi'(a)=c\ne0\) 及所需的幂关系。

对 \(\phi\) 应用\cref{b3:thm:01-local-inverse}，取得 \(a\) 的开邻域 \(U_0\) 和 \(0\) 的开邻域 \(V_0\)，使 \(\phi:U_0\to V_0\) 及其逆函数 \(\psi\) 均全纯，且 \(\phi'\) 在 \(U_0\) 无零。选 \(r>0\)，使 \(D(0,r)\subseteq V_0\)，再令
\[
 U=\{z\in U_0:\phi(z)\in D(0,r)\}.
\]
连续性使 \(U\) 开，且它包含 \(a\)；限制后的 \(\phi\) 把 \(U\) 双射到整个 \(D(0,r)\)，逆函数为 \(\psi\) 在此圆盘上的限制。这样取得的源邻域 \(U\) 未必是几何圆盘，但它在 \(\phi\) 坐标中恰是一个圆盘，幂关系在其上成立。
\end{proof}

当 \(m=1\) 时，原函数本身已在邻域内可逆；当 \(m\geq2\) 时，新坐标中的函数为 \(\xi\mapsto f(a)+\xi^m\)。构造中开根只针对始终靠近非零值的局部因子，没有要求含零点的函数存在全纯根，也没有要求任意区域上的非零函数都有全纯对数。

\Needspace{6\baselineskip}
\begin{theorem}[开映射]\label{b3:thm:01-open-mapping}
设 \(\Omega\subseteq\mathbb C\) 为区域，\(f:\Omega\to\mathbb C\) 为非常值全纯函数。则 \(f\) 是\term[kai-yingshe]{开映射}[open map]：对任意开集 \(O\subseteq\Omega\)，像 \(f(O)\) 在 \(\mathbb C\) 中开。这个结论称为\term[kai-yingshe-dingli]{开映射定理}[open mapping theorem]。
\end{theorem}
\begin{proof}
空集的像为空集。若 \(O\ne\varnothing\)，任取 \(a\in O\)。恒等定理说明 \(f-f(a)\) 不会在 \(a\) 附近恒零，否则 \(f\) 将在整个 \(\Omega\) 上为常数。将\cref{b3:prop:01-local-normal-form}的邻域取在 \(O\) 内，得到 \(f(z)=f(a)+\phi(z)^m\)，且 \(\phi(U)=D(0,r)\)。幂映射把 \(D(0,r)\) 满射到 \(D(0,r^m)\)：任意非零数 \(w\) 的 \(m\) 次根均有模 \(|w|^{1/m}\)，零的原像为零。因此
\[
D\bigl(f(a),r^m\bigr)=f(U)\subseteq f(O).
\]
像中每个点都是内点，故 \(f(O)\) 开。
\end{proof}

\Needspace{6\baselineskip}
\begin{corollary}\label{b3:cor:01-local-injectivity}
设 \(\Omega\subseteq\mathbb C\) 为开集，\(a\in\Omega\)，\(f:\Omega\to\mathbb C\) 全纯。则 \(f\) 在 \(a\) 的某邻域内单射，当且仅当 \(f'(a)\ne0\)。更一般地，若 \(f-f(a)\) 在 \(a\) 具有有限零点阶数 \(m\ge1\)，\(U\subseteq\Omega\) 是 \(a\) 的开邻域，且全纯双射 \(\phi:U\to D(0,r)\)、\(r>0\)，满足 \(\phi(a)=0\)、\(f=f(a)+\phi^m\)，则每个满足 \(0<|w-f(a)|<r^m\) 的值 \(w\) 恰有 \(m\) 个位于 \(U\) 的不同原像。
\end{corollary}
\begin{proof}
先数标准形邻域内的原像。写 \(w-f(a)=\rho e^{\mathrm{i}\theta}\)，其中 \(0<\rho<r^m\)。方程 \(f(z)=w\) 等价于 \(\phi(z)^m=w-f(a)\)，其 \(m\) 个根为
\[
 \xi_j=\rho^{1/m}\exp\left(\frac{\mathrm{i}(\theta+2\pi j)}m\right),
 \quad j=0,\ldots,m-1.
\]
这些根两两不同，均落在 \(D(0,r)\)，且由模和辐角的表示不存在其他根。\(\phi\) 的双射性给出恰好 \(m\) 个不同原像。

若 \(f'(a)=0\) 而 \(f\) 不局部恒定，则其首个非零 Taylor 项的次数 \(m\ge2\)。对任意给定的 \(a\) 的开邻域，都可在其中取标准形邻域，再选择 \(0<|w-f(a)|<r^m\)，得到多个原像，所以该邻域上不能单射；局部恒定时也不能在任何 \(a\) 的邻域内单射。反方向由局部逆函数定理得到。
\end{proof}

函数 \(z\mapsto z^m\) 在 \(m\ge2\) 时提供最直接的例子：原点处导数为零，但每个小圆盘仍映为一个圆盘。开映射要求像含有邻域，局部单射要求不同点不重合。对实可微函数 \(z\mapsto|z|^2\)，小圆盘的像只落在实轴上，开映射结论便不成立。

\ProseParagraphBreak

上述原像计数还可以写成具体的局部逆分支。令 \(b=f(a)\)，\(\psi=\phi^{-1}:D(0,r)\to U\)，并在目标圆盘中删去由 \(b\) 向负实方向伸出的半径，得到
\[
 W=\{w:0<|w-b|<r^m,\ w-b\notin(-\infty,0]\}.
\]
对 \(j=0,\ldots,m-1\)，定义
\[
 h_j(w)=\psi\left(e^{2\pi\mathrm{i}j/m}\cdot
          \exp\bigl(\tfrac1m\operatorname{Log}(w-b)\bigr)\right),
 \quad w\in W.
\]
括号中数的模为 \(|w-b|^{1/m}<r\)，所以它在 \(\psi\) 的定义域内。各 \(h_j\) 全纯，满足 \(f(h_j(w))=w\)，并且对每个 \(w\in W\) 恰好给出 \(U\) 内全部原像。这里删去半径是为了在同一目标开集上连续标定这些根；对单个非零目标值计数时，不需要作任何分支选择。

\ProseParagraphBreak

对开集上的全纯函数，导数为零的点称为\term[linjie-dian]{临界点}[critical point]，其像称为\term[linjie-zhi]{临界值}[critical value]。在上述标准形中，\(\phi'\) 处处非零，而
\[
 f'(z)=m\phi(z)^{m-1}\phi'(z).
\]
故 \(m\ge2\) 时，\(a\) 是 \(U\) 内唯一的临界点，\(b\) 是这部分映射的唯一临界值；对 \(w\ne b\)，那些不同原像处的导数均不为零，每一个原像附近都有通常的局部逆。其他源邻域是否也映到这些值，需要另行研究，不能由这一局部计数决定。

\ProseParagraphBreak

具体来看 \(f(z)=z^2+z^3=z^2(1+z)\)。在 \(|z|<1/2\) 上，\(1+z\) 落在右半平面，可以定义
\[
 \phi(z)=z\exp\bigl(\tfrac12\operatorname{Log}(1+z)\bigr).
\]
直接求值得 \(\phi(0)=0\)、\(\phi'(0)=1\)，并且 \(f=\phi^2\)。按局部标准形的做法，进一步限制到 \(\phi:U\to D(0,r)\) 可逆的邻域，记逆函数为 \(\psi\)。在 \(0<|w|<r^2\) 且删去负实半径的目标区域中，两条逆分支正是
\[
 z_+(w)=\psi\bigl(\sqrt w\bigr),
 \quad z_-(w)=\psi\bigl(-\sqrt w\bigr),
 \quad \sqrt w=\exp\bigl(\tfrac12\operatorname{Log}w\bigr).
\]
三次项通过坐标 \(\phi\) 保留下来。它会改变根的位置，却不改变足够小的 \(U\) 内恰有两个不同原像这一事实。

\ProseParagraphBreak

还可以算出位置修正。由 Taylor 展开，\(\phi(z)=z+z^2/2+O(|z|^3)\)。因 \(\psi'(0)=1\)，写 \(\psi(\xi)=\xi+A\xi^2+O(|\xi|^3)\)，代入 \(\phi(\psi(\xi))=\xi\)，比较二次项得 \(A=-1/2\)。因此
\[
\begin{aligned}
 z_+(w)&=\sqrt w-\tfrac12w+O(|w|^{3/2}),\\
 z_-(w)&=-\sqrt w-\tfrac12w+O(|w|^{3/2}).
\end{aligned}
\]
余项估计来自 \(\psi\) 在一个更小闭圆盘上的 Taylor 余项，所以在所取割圆盘中对所有趋近方向同时有效。若沿 \(w(t)=\rho e^{\mathrm{i}t}\)、\(0\le t\le2\pi\)，连续选根从 \(\sqrt\rho\) 出发，便得到 \(\sqrt\rho e^{\mathrm{i}t/2}\)：任意连续选根与这个表达式的比只能取 \(1\) 或 \(-1\)，连续性使该比保持常数。终点为 \(-\sqrt\rho\)，两条局部反解交换，因而不能各自成为整个去心目标圆盘上的连续单值逆函数。

\ProseParagraphBreak

本例中 \(f'(z)=z(2+3z)\)，两个临界点为 \(0\)、\(-2/3\)，对应临界值为 \(0\)、\(4/27\)。在第二个点令 \(z=-2/3+\zeta\)，则
\[
 f(-2/3+\zeta)=\frac4{27}-\zeta^2+\zeta^3,
\]
所以那里同样是二次的局部模型。另一方面，\(f(-1)=0\)、\(f'(-1)=1\)，局部逆函数定理还给出靠近 \(-1\) 的一个分支 \(z_0(w)=-1+w+O(|w|^2)\)。它与靠近原点的两个分支位于不交源邻域。由此，充分小的非零 \(w\) 在全平面有三个不同原像，而在原点的标准形邻域内只有两个；次数为三的多项式至多有三个不同根，因而这里已经列尽全部原像。局部重数记录的是原像在哪个点聚合，不是整个定义域中的总原像数。

\ProseParagraphBreak

\chapref{b3:ch:02}将讨论一般曲线上的复积分，以及 Cauchy 积分公式给出的估计与全局结论。开映射也将用来证明内部最大模原理。本章得到的局部幂级数、零点分解和标准形，随后还会用于奇点分类、零点计数与解析延拓。

% 本章小结（不计入编号）
\begin{chaptersummary}
复可微把二维实微分限制为复数乘法。偏导数满足 Cauchy--Riemann 方程只是其中一部分；若没有实可微的统一余项控制，复差商仍可能不存在。反过来，函数在开集上处处复可微，则由局部积分公式得到收敛的 Taylor 展开，各阶导数的存在与连续性随之成立。

\ProseParagraphBreak

幂级数逐项运算的依据是收敛圆内部的统一估计。选取两个有间隔的半径，可用几何衰减同时控制尾项和各阶导数；把内半径推到外半径时，原有估计可能失效，因此边界需要另行讨论。指数与三角函数可在全平面构造，对数却需要选定分支。局部存在全纯的对数或逆函数，并不保证能在整个定义域连续选取；指数函数的周期和去心平面上不存在连续单值对数的事实说明了这一区别。

\ProseParagraphBreak

非恒零全纯函数的零点由最先出现的非零 Taylor 项控制。恒等定理通过连通性把局部相等传到整个区域；局部标准形进一步把非常值函数化成一个幂映射。标准坐标中，导数非零对应一次幂，较高阶零点对应多次幂；附近非中心值的原像数目便由这个次数确定。这里计数的是选定邻域内的原像，其他位置仍可能另有逆分支。两种情形都把邻域映为邻域。下一章将在本章局部积分结果上继续建立一般 Cauchy 理论。
\end{chaptersummary}

\BookExercises

\Cref{b3:ex:01-reconstruct-real-part}与\cref{b3:ex:01-power-series-ode}训练方程和幂级数的具体计算；\cref{b3:ex:01-mean-value-failure}检验实变量中值定理的适用边界。\Cref{b3:ex:01-disk-distance-ratio}考察圆盘变换的展开与参数极限，\cref{b3:ex:01-equal-modulus}与\cref{b3:ex:01-local-involution}组合使用零点、恒等定理和局部求逆。两道带星号的题进一步讨论逆函数的展开半径和幂级数跨越收敛圆周的障碍，所需工具均已在本章建立。

% 自拟题组；保留既有训练方向，解答使用本章正文工具，不复刻第三方原题。
% 实质增量：实部重建、中值等式反例、整函数常微分方程、圆盘变换的参数极限、
% 等模函数共同零点、局部对合线性化、逆函数半径与阶乘缺项的边界障碍。

\Needspace{6\baselineskip}
\begin{exercise}\label{b3:ex:01-reconstruct-real-part}
求出所有满足
\[
  \operatorname{Re}f(x+\mathrm{i}y)=x^3-3xy^2+2xy
  \quad(x,y\in\mathbb R)
\]
的整函数 \(f:\mathbb C\rightarrow\mathbb C\)，并说明给定条件为什么只留下一个实常数的自由度。
\end{exercise}
\begin{solution}
记题给的实部为 \(u\)，虚部为 \(v\)。全纯性保证实部、虚部的偏导数存在并满足 Cauchy--Riemann 方程，因而
\[
  v_y=u_x=3x^2-3y^2+2y,
  \quad v_x=-u_y=6xy-2x.
\]
固定 \(x\)，对第一个等式关于 \(y\) 积分，得到
\[
  v(x,y)=3x^2y-y^3+y^2+h(x),
  \quad h(x)=v(x,0).
\]
这里的积分常数可以依赖 \(x\)。由于 \(h\) 可微，再代入第二个等式，有
\(6xy+h'(x)=6xy-2x\)，所以 \(h(x)=-x^2+C\)，其中 \(C\in\mathbb R\)。因此
\[
  f(z)=z^3-\mathrm{i}z^2+\mathrm{i}C,
  \quad C\in\mathbb R.
\]
反向核对时，\(\operatorname{Re}(z^3)=x^3-3xy^2\)、\(\operatorname{Re}(-\mathrm{i}z^2)=2xy\)，所以这些整函数都符合题设。计算已从必要条件确定了全部可能的虚部，故没有遗漏。固定实部后只能给虚部加一个实常数；若给 \(f\) 加上实部非零的常数，就会改变题给实部。
\end{solution}

\Needspace{6\baselineskip}
\begin{exercise}\label{b3:ex:01-mean-value-failure}
设 \(U\subseteq\mathbb C\) 为开集，\(f:U\rightarrow\mathbb C\) 全纯，\(a,b\in U\)、\(a\ne b\)，且整条闭线段 \([a,b]\) 包含于 \(U\)。证明
\[
  f(b)-f(a)=(b-a)\int_0^1 f'\bigl(a+t(b-a)\bigr)\,\mathrm{d}t
\]
及相应的模估计。再构造一个整函数，说明一般不存在 \(t\in(0,1)\) 使
\[
  f(b)-f(a)=(b-a)\cdot f'\bigl(a+t(b-a)\bigr).
\]
\end{exercise}
\begin{solution}
由全纯与解析的等价性，\(f'\) 连续。函数
\(F(t)=f\bigl(a+t(b-a)\bigr)\) 在 \([0,1]\) 上是复值 \(C^1\) 函数，并满足
\[
  F'(t)=(b-a)\cdot f'\bigl(a+t(b-a)\bigr).
\]
将实变量微积分基本定理分别用于 \(F\) 的实部和虚部，合并后便得到题给的积分等式。三角不等式进一步给出
\[
\begin{aligned}
  |f(b)-f(a)|
  &\le |b-a|\cdot
       \int_0^1\left|f'\bigl(a+t(b-a)\bigr)\right|\,\mathrm{d}t\\
  &\le |b-a|\cdot
       \max_{0\le t\le1}\left|f'\bigl(a+t(b-a)\bigr)\right|.
\end{aligned}
\]
最后一个最大值存在，因为被取模的函数在紧区间上连续。

取 \(f(z)=e^z\)、\(a=0\)、\(b=2\pi\mathrm{i}\)。这时 \(f(b)-f(a)=0\)，而
\[
  (b-a)\cdot f'\bigl(a+t(b-a)\bigr)
  =2\pi\mathrm{i}e^{2\pi\mathrm{i}t}\ne0
\]
对所有实数 \(t\) 都成立。因此，实变量中值定理中的单点等式不能照搬。复导数沿线段的方向可以相互抵消，使积分为零，而没有任何一个导数值等于零。
\end{solution}

\Needspace{6\baselineskip}
\begin{exercise}\label{b3:ex:01-power-series-ode}
求所有满足
\[
  f'(z)=z\cdot f(z)\quad(z\in\mathbb C),
  \quad f(0)=1
\]
的整函数 \(f:\mathbb C\rightarrow\mathbb C\)。先由幂级数求出系数并核对收敛半径，再证明解的唯一性。
\end{exercise}
\begin{solution}
任何可能的解在原点附近都具有展开 \(f(z)=\sum_{n\ge0}a_nz^n\)。逐项求导后比较同次幂系数，得到
\[
  a_0=1,\quad a_1=0,
  \quad (n+1)a_{n+1}=a_{n-1}\quad(n\ge1).
\]
这个递推把奇数次与偶数次系数分开。所有奇数次系数均为零，而
\[
  a_{2k}=\frac{a_{2k-2}}{2k}
         =\frac1{2^k k!}\quad(k\ge1).
\]
因此唯一的候选级数为
\[
  f(z)=\sum_{k=0}^{\infty}\frac{z^{2k}}{2^k k!}
       =e^{z^2/2}.
\]
对任意固定的 \(r>0\)，在 \(|z|\le r\) 上，绝对值级数由
\(\sum_{k\ge0}r^{2k}/(2^k k!)\) 控制。其相邻项比值为
\(r^2/(2(k+1))\)，趋于零，故控制级数收敛。所求幂级数在每个闭圆盘上一致绝对收敛，半径为无穷大；逐项求导和上述递推证明其确实满足微分方程与初值。

若 \(g\) 是另一整函数解，则同样的递推使它与 \(f\) 在原点附近的全部 Taylor 系数相同。因此它们在一个开圆盘内相等，再由恒等定理在连通的 \(\mathbb C\) 上相等。也可以直接计算积分因子：
\[
  \bigl(e^{-z^2/2}\cdot g(z)\bigr)'
  =e^{-z^2/2}\cdot\bigl(g'(z)-z\cdot g(z)\bigr)=0.
\]
其值于是局部恒定，再由连通性为同一个常数；原点处的值为 \(1\)，再次得到 \(g(z)=e^{z^2/2}\)。
\end{solution}

\Needspace{6\baselineskip}
\begin{exercise}\label{b3:ex:01-disk-distance-ratio}
设 \(a\in\mathbb C\)、\(0<|a|<1\)，在 \(\mathbb C\setminus\{1/\overline a\}\) 上定义
\[
  \phi_a(z)=\frac{z-a}{1-\overline a z}.
\]
求它在原点的 Taylor 级数、精确收敛半径，以及截到 \(z^N\) 后在 \(|z|\le r\) 上的余项界，其中 \(N\ge0\) 为整数，\(0<r<1/|a|\)。再固定 \(\zeta\in\mathbb C\)、\(|\zeta|=1\)，令 \(a=t\zeta\)，其中 \(0<t<1\)。证明当 \(t\uparrow1\) 时，\(\phi_{t\zeta}\) 在每个固定的闭圆盘 \(\overline{D(0,r)}\)、\(0<r<1\) 上一致趋于常数 \(-\zeta\)，却不在闭单位圆盘上一致趋于该常数。
\end{exercise}
\begin{solution}
先把常数项分离：
\[
  \phi_a(z)=-a+(1-|a|^2)\frac{z}{1-\overline a z}.
\]
在 \(|z|<1/|a|\) 上应用几何级数，得
\[
  \phi_a(z)=-a+(1-|a|^2)
       \sum_{n=1}^{\infty}\overline a^{\,n-1}z^n.
\]
对 \(n\ge1\)，系数模的 \(n\) 次根为
\[
  (1-|a|^2)^{1/n}|a|^{(n-1)/n}\longrightarrow|a|.
\]
所以 Cauchy--Hadamard 公式给出的半径恰为 \(1/|a|\)。由于 \(a\ne0\) 且 \(|a|<1\)，这些高次系数并不消失。

设 \(P_N\) 为保留到 \(z^N\) 的多项式，其中 \(P_0=-a\)。对余下的几何级数求和，得到
\[
  \phi_a(z)-P_N(z)
  =\frac{(1-|a|^2)\overline a^{\,N}z^{N+1}}
         {1-\overline a z}.
\]
在 \(|z|\le r<1/|a|\) 上有
\(|1-\overline a z|\ge1-|a|\cdot|z|\ge1-|a|r>0\)，因而
\[
  \sup_{|z|\le r}|\phi_a(z)-P_N(z)|
  \le\frac{(1-|a|^2)|a|^N r^{N+1}}{1-|a|r}.
\]
这个界对 \(N=0\) 同样成立。

最后直接计算
\[
  \phi_{t\zeta}(z)+\zeta
  =\frac{(1-t)(z+\zeta)}{1-t\overline\zeta z}.
\]
在固定的 \(|z|\le r<1\) 上，其模不超过
\((1-t)(1+r)/(1-tr)\)。分子趋于零，分母至少为 \(1-r>0\)，所以得到所需一致收敛。对于每个 \(t<1\)，分母在闭单位圆盘上也不为零，故边界点值均有定义；但
\[
  \phi_{t\zeta}(\zeta)=\frac{(1-t)\zeta}{1-t}=\zeta,
  \quad |\phi_{t\zeta}(\zeta)+\zeta|=2.
\]
于是闭单位圆盘上的一致收敛失败。固定内圆盘的正边界间隔正是前一估计能够统一控制分母的原因。
\end{solution}

\Needspace{6\baselineskip}
\begin{exercise}\label{b3:ex:01-equal-modulus}
设 \(U\subseteq\mathbb C\) 为非空连通开集，\(f,g:U\rightarrow\mathbb C\) 全纯，且
\(|f(z)|=|g(z)|\) 对每个 \(z\in U\) 成立。证明存在常数 \(c\in\mathbb C\)、\(|c|=1\)，使 \(f=cg\) 于 \(U\)。解答须处理共同零点，不能直接假设 \(f/g\) 处处有定义。
\end{exercise}
\begin{solution}
若 \(g\equiv0\)，等模条件使 \(f\equiv0\)，取 \(c=1\) 即可。以下设 \(g\not\equiv0\)，因而 \(f\not\equiv0\)。它们的零点集合相同，且由恒等定理和零点分解，每个共同零点都有有限阶。

任取一个共同零点 \(a\)。在它的某个圆盘邻域内写
\[
  f(z)=(z-a)^m\cdot F(z),
  \quad g(z)=(z-a)^n\cdot G(z),
  \quad F(a)\ne0,\quad G(a)\ne0.
\]
这里 \(m,n\) 为正整数，\(F,G\) 全纯。缩小该圆盘，使 \(F,G\) 均无零点。对其中的 \(z\ne a\)，等模条件化为
\[
  |z-a|^{m-n}\cdot\left|\frac{F(z)}{G(z)}\right|=1.
\]
当 \(z\to a\) 时，商的模趋于有限正数 \(|F(a)/G(a)|\)。如果 \(m>n\)，左边趋于零；如果 \(m<n\)，左边趋于无穷大。因此只能有 \(m=n\)。这时在去心圆盘上 \(f/g=F/G\)，而 \(F/G\) 已在整个圆盘上全纯，故可用其在 \(a\) 的值补定义商。

对每个共同零点作这样的补定义，并在 \(g\ne0\) 处仍取通常的商，得到 \(U\) 上的函数 \(h\)。它在每点附近都等于一个全纯函数，所以全纯；各局部定义在交集上相容，因为它们在非零点处都等于 \(f/g\)，在零点处的值再由连续性唯一确定。由非零点的等模条件以及零点处的连续性，\(|h|=1\) 在整个 \(U\) 上成立。

若 \(h\) 非常值，开映射定理将使 \(h(U)\) 成为非空开集；但 \(h(U)\) 包含于单位圆，后者不含任何开圆盘。这是矛盾。因此 \(h=c\)，其中 \(|c|=1\)。在非零点处有 \(f=cg\)，共同零点处两边均为零，故所求等式在全域成立。
\end{solution}

\Needspace{6\baselineskip}
\begin{exercise}\label{b3:ex:01-local-involution}
设 \(U\subseteq\mathbb C\) 是包含原点的区域，\(f:U\rightarrow U\) 全纯，并满足
\[
  f(0)=0,\quad f\bigl(f(z)\bigr)=z\quad(z\in U).
\]
证明 \(f'(0)\) 只能为 \(1\) 或 \(-1\)。若为 \(1\)，证明 \(f(z)=z\) 于 \(U\)；若为 \(-1\)，构造一个在原点附近具有全纯逆的函数 \(h\)，使
\(h\bigl(f(z)\bigr)=-h(z)\)，并说明如何选取原点邻域，使这一等式确实把 \(f\) 化成新坐标中的 \(w\mapsto-w\)。
\end{exercise}
\begin{solution}
在原点对复合恒等式使用链式法则，得到
\(f'(f(0))\cdot f'(0)=1\)，即 \(f'(0)^2=1\)。

先设 \(f'(0)=1\)。若 \(f\) 不恒等于恒等映射，则恒等定理排除了 \(f(z)-z\) 在原点附近恒零的可能。由 Taylor 展开，存在最小的整数 \(m\ge2\) 及 \(a\ne0\)，使
\[
  f(z)=z+az^m+R(z),\quad R(z)=O(|z|^{m+1}).
\]
为了核对复合的首个非零项，先注意 \(f(z)=O(|z|)\)，故
\(R(f(z))=O(|z|^{m+1})\)。又由
\(f(z)=z\bigl(1+O(|z|^{m-1})\bigr)\) 得
\[
  f(z)^m=z^m+O(|z|^{2m-1})
         =z^m+O(|z|^{m+1}),
\]
最后一步用了 \(m\ge2\)。因此
\[
\begin{aligned}
  f\bigl(f(z)\bigr)
  &=f(z)+a\cdot f(z)^m+R(f(z))\\
  &=z+2az^m+O(|z|^{m+1}).
\end{aligned}
\]
它不可能恒等于 \(z\)：否则减去 \(z\) 后除以 \(z^m\)，令 \(z\to0\)，便有 \(2a=0\)。这个矛盾证明 \(f(z)=z\) 于整个 \(U\)。

再设 \(f'(0)=-1\)。在 \(U\) 上令
\[
  h(z)=\frac{z-f(z)}2.
\]
则 \(h(0)=0\)、\(h'(0)=1\)，且对所有 \(z\in U\)，
\[
  h\bigl(f(z)\bigr)
  =\frac{f(z)-f(f(z))}2
  =-h(z).
\]
由局部逆函数定理，可取原点邻域 \(V_0\subseteq U\)，使 \(h\) 在 \(V_0\) 上为具有全纯逆的双射。还须保证所选邻域在 \(f\) 下保持不变。题设 \(f\circ f=\operatorname{id}\) 说明 \(f\) 是自己的连续逆，故
\[
  V_1=V_0\cap f(V_0)
\]
是包含原点的开集，并满足 \(f(V_1)=V_1\)。限制后的 \(h:V_1\rightarrow h(V_1)\) 仍为全纯双射，像是原点的开邻域。取 \(r>0\) 使 \(D(0,r)\subseteq h(V_1)\)，并令
\[
  V=(h|_{V_1})^{-1}(D(0,r)).
\]
若 \(z\in V\)，则 \(f(z)\in V_1\)，并由 \(h(f(z))=-h(z)\in D(0,r)\) 得 \(f(z)\in V\)。再次应用 \(f\) 可知 \(f(V)=V\)。因此 \(h:V\rightarrow D(0,r)\) 确实给出同一邻域内的坐标变换，并且
\[
  (h\circ f\circ h^{-1})(w)=-w
  \quad(|w|<r).
\]
\end{solution}

\Needspace{6\baselineskip}
\begin{optionalexercise}\label{b3:ex:01-inverse-radius}
设 \(F:\mathbb C\rightarrow\mathbb C\)、\(F(z)=z+z^2\)。求满足 \(G(0)=0\) 的全纯局部逆 \(G\) 在原点的前四个非零 Taylor 系数，并证明该 Taylor 级数的收敛半径恰为 \(1/4\)。这里 \(G\) 的初始定义域只要求是原点的一个开邻域。
\end{optionalexercise}
\begin{solution}
因为 \(F'(0)=1\)，局部逆函数定理保证所需 \(G\) 存在；限制到充分小的邻域后，它也由 \(G(0)=0\) 唯一确定。写
\[
  G(w)=b_1w+b_2w^2+b_3w^3+b_4w^4+O(|w|^5).
\]
代入 \(G+G^2=w\)，比较前四次幂的系数，依次得到
\[
\begin{aligned}
  b_1&=1,&b_2+b_1^2&=0,\\
  b_3+2b_1b_2&=0,&b_4+2b_1b_3+b_2^2&=0.
\end{aligned}
\]
因此
\[
  G(w)=w-w^2+2w^3-5w^4+O(|w|^5).
\]

接着证明收敛半径的下界。若 \(|w|<1/4\)，则
\(\operatorname{Re}(1+4w)>0\)，主对数有定义。函数
\[
  H(w)=\frac{\exp\bigl(\tfrac12\operatorname{Log}(1+4w)\bigr)-1}{2}
\]
在整个 \(D(0,1/4)\) 上全纯。令分子中的指数值为 \(S(w)\)，则
\(S(w)^2=1+4w\)、\(S(0)=1\)，从而 \(H(0)=0\) 且
\[
  H(w)+H(w)^2=\frac{S(w)^2-1}{4}=w.
\]
在充分小的原点邻域内，\(H(w)\) 落在 \(F\) 的局部单射邻域，故 \(F(H(w))=w\) 迫使 \(H=G\)。于是 \(H\) 的原点 Taylor 系数就是所求局部逆的系数。由圆盘上的全纯函数具有覆盖整个圆盘的 Taylor 展开，所求级数的收敛半径 \(R\ge1/4\)。

若 \(R>1/4\)，记该 Taylor 级数在 \(D(0,R)\) 上的和为 \(\widetilde G\)。它在原点附近与 \(G\) 相同，故
\(\widetilde G+\widetilde G^2-w\) 在那里为零。这一表达式在连通圆盘 \(D(0,R)\) 上全纯，恒等定理便使
\[
  \widetilde G(w)+\widetilde G(w)^2=w
  \quad(|w|<R)
\]
在全盘成立，而不只在最初的逆函数邻域成立。特别地，\(w_0=-1/4\) 属于此圆盘，代入得到
\((\widetilde G(w_0)+1/2)^2=0\)，所以 \(\widetilde G(w_0)=-1/2\)。但对全盘恒等式求导有
\[
  \bigl(1+2\widetilde G(w)\bigr)\cdot\widetilde G'(w)=1.
\]
在 \(w=w_0\) 处左边为零，矛盾。因此 \(R=1/4\)。这里出现的障碍来自 \(F'(-1/2)=0\)，而 \(F(-1/2)=-1/4\) 正好位于所求收敛圆周上。
\end{solution}

\Needspace{6\baselineskip}
\begin{optionalexercise}\label{b3:ex:01-factorial-gaps}
在单位圆盘 \(D(0,1)\) 上考虑
\[
  f(z)=\sum_{n=1}^{\infty}z^{n!}.
\]
证明这确实给出一个全纯函数，且其幂级数的收敛半径为 \(1\)。再证明：对每个 \(\zeta\in\mathbb C\)、\(|\zeta|=1\)，都不存在 \(\varepsilon>0\) 及 \(V=D(\zeta,\varepsilon)\) 上的全纯函数 \(F:V\rightarrow\mathbb C\)，使 \(F=f\) 在 \(V\cap D(0,1)\) 上成立。
\end{optionalexercise}
\begin{solution}
把缺少的幂次系数补成零，可写成 \(f(z)=\sum_{k=0}^{\infty}c_kz^k\)。由于 \(1!,2!,3!,\ldots\) 严格递增，系数 \(c_k\) 在 \(k=n!\) 时为 \(1\)，其余为 \(0\)。所以
\[
  \limsup_{k\to\infty}|c_k|^{1/k}=1.
\]
Cauchy--Hadamard 公式给出半径 \(1\)，并保证在每个闭子圆盘上一致绝对收敛；由幂级数逐项求导定理，其和在单位圆盘上全纯。

先考虑一个单位根 \(\eta\)，即存在正整数 \(q\) 使 \(\eta^q=1\)。取整数 \(N\ge q\)。当 \(n\ge N\) 时，\(q\) 整除 \(n!\)，所以 \(\eta^{n!}=1\)。对 \(0<r<1\)，有
\[
  f(r\eta)
  =\sum_{n<N}(r\eta)^{n!}+T(r),
  \quad T(r)=\sum_{n\ge N}r^{n!}.
\]
前一有限和的模不超过 \(N-1\)。后一和的每项都是非负实数；任意固定正整数 \(L\) 后，
\[
  T(r)\ge\sum_{n=N}^{N+L-1}r^{n!}
       \longrightarrow L\quad(r\uparrow1).
\]
因此 \(\liminf_{r\uparrow1}T(r)\ge L\) 对每个 \(L\) 成立，故 \(T(r)\to+\infty\)。于是
\[
  \operatorname{Re}f(r\eta)\ge T(r)-(N-1)\longrightarrow+\infty.
\]
这说明 \(f\) 沿趋向任意单位根的半径都没有有限极限，因而不可能等于该单位根邻域内某个连续函数的限制，更不可能全纯延拓到那里。

单位根 \(e^{2\pi\mathrm{i}p/q}\) 在单位圆上稠密：这是有理数在实轴上的稠密性与映射 \(t\mapsto e^{2\pi\mathrm{i}t}\) 的连续性给出的。若题设的 \(V,F\) 对某个 \(\zeta\) 存在，可以在开圆盘 \(V\) 中选取一个单位根 \(\eta\)。当 \(r<1\) 足够接近 \(1\) 时，\(r\eta\in V\cap D(0,1)\)，故
\[
  f(r\eta)=F(r\eta)\longrightarrow F(\eta).
\]
右边有限，与刚得到的实部趋于无穷大矛盾。由此排除了跨越单位圆上任意一点的所述局部延拓。
\end{solution}
